\section{Opening stanzas}
\begin{verse}
	\labelh{OSTZ1}
  In\refm{2}%\pratika{yasmin na}
\pratikap{p01_1}\labelh{FVFPTRL} Him \karman\ and its ripening do not exist,%
  \footnote{\label{NONFSFP}This is a reference to \rYS{1.24}{1.24}: \emph{kleśakarmavipākāśayair aparāmṛṣṭaḥ}.  Impurities (\klesa\/s) are mentioned in the next pāda.%
  }
  From Him they originated,%
  \footnote{%The verb \emph{āstām} may be read in several ways: as the third person dual, imperfect of the root \emph{as-} (2nd class) in the sense ``the two (\karman\/s and their ripening) originated from/because of Him''; if we read the phrase as \emph{yataḥ} (ablative) + \emph{ā} + \emph{astām}, even the interpretation, ``may the two exist up to Him (viz., function except for Him),'' is possible.  %
	The part \emph{[karmavipākau] yata āstām} may be interpreted in two ways.  One way is to understand it to mean that the teachings on \karman\/s and their results, i.e., the \Veda s, were given by \Isvara.  That \karman\/s and their ripening, in the context of \dharma, or in the sense of Vedic rituals and their fruition, were taught by \Isvara\ is expressed a few times in this text.  See syllogisms following the one starting with \emph{anekakartṛ\/°} \reftoeditionpar{p25_39} translated on p.~\pageref{VRNASRM25}.  Also, in one interpretation of the part of \rYBh{1.24}{1.24} where it asserts that \Isvara\ is the source (\nimitta) of \sastra\ and that \sastra\ is the source of \Isvara\ as well (starting with the paragraph \emph{atro\-cyate} \reftoeditionpar{p24_13} translated on p.~\pageref{SSTRISVR24}), the author interprets the word \nimitta\ to mean the cause.  The \sastra\ in this case most likely includes the Vedas.  Another indication that the author had the Vedas in mind in the expression \emph{[karmavipākau] yata āstām} is his marked hostility towards the \Mimamsaka s. \Sabara's interpretation of the word \atha\ is examined shortly (see pp.~\pageref{SBRATH}\,ff.);  he alludes to \Kumarila\ many times as an opponent in the \Isvara\ section. \Kumarila\ champions the view that the Vedas have no personal author and strongly skeptical of the creator god.  Our author even cites and criticizes \Kumarila's view on \sphota\ in the commentary on \rYS{3.17}{3.17}. It is natural for such an author to mention the relationship between \Isvara\ and the Vedas in the opening stanza.%

	Another likely meaning of \emph{[karmavipākau] yata āstām} is that the rules regarding \karman\/s (in the sense of deeds) and their ripening were set by \Isvara\ and he oversees their operation.  This old idea is mentioned several times in the \BhGBh\ in the form of \emph{sarvasya karma\-phala\-jātasya dhātāraṃ vidhātāraṃ} (\rBhGBh{08.9}{8.9}), \emph{dhātā karma\-phalasya prāṇibhyo vidhātā} (\rBhGBh{09.17}{9.17}), \emph{dhātāhaṃ karma\-phalasya vidhātā} (\rBhGBh{10.33}{10.33}), etc.  Since the author shows affiliation to the \BhG\ in the third pāda of this stanza by using the verb \kal\ in relation to \kala, it is conceivable that he had this old idea in mind.  Also, one area of \Kumarila's criticism of \Isvara\ involves the relationship between the law of \karman\/s (or \dharma) and \Isvara\ (\SVSAP{043b, 50, 69--72, 75}{kk.\ 43b, 50, 69--72, 75}).  One common theme that appears in \SVSAP{072 and 75}{kk.\ 72 and 75} is that the involvement of \Isvara\ in the workings of \karman\/s and their fruits is redundant.  It is also conceivable that the author wanted to give an answer to this problem.%
	}\\
  Impurities (\klesa\/s) are not enough [to affect] him; but not surmountable for everybody [else], \\%
  He cannot be fixed by the compelling eye of \inidx{Kāla},%
  \footnote{The readings in two primary manuscripts differ. \Tm\ reads \emph{kāladiśā} and L reads \emph{kāladṛśā}.  Both seem possible.  Even the singular for the dual \DD\ compound \emph{kāladiś} may still be acceptable because its etymologically related synonym \emph{deśakāla} is allowed to be used as singular.  However, I slightly favor the reading \emph{kāladṛśā} because it is less likely for it to change to \emph{kāladiśā} (more familiar sounding and graphically simpler) than the other way around.  This part refers to \rYS{1.26}{1.26}: \emph{sa pūrveṣām api guruḥ kālenānavacchedāt}.%
  } \\
  I salute Him, the Lord of the world, the enemy of \Kaitabha.%
  \footnote{\label{VSNSM1}This stanza pays homage to \Visnu\ (the enemy, i.e., the slayer, of \Kaitabha).  On our author's affiliation to the \Vaisnavism, see pp.\ \pageref{NRYNN25}, \pageref{NNRYN26}, \pageref{VISNU27}, etc.  Despite the common view that associates \Sankara\ with \Saivism, it has been pointed out that \Sankara\ had a religious background closer to \Vaisnavism.  See \citet{hacker:relation}.  Cf.\ \citet[4]{mayeda:thousand}.}%  On the other hand, however, we may practice some caution assuming the genuinness of the verse since we have evidence not all the closing verses are genuine.  See p.~\pageref{CLSVRS} in the introduction.}
  
  He\refm{3}%\pratika{yas sarvavit}
\pratikap{p01_2} is omniscient, ruling all, and omnipotent,\footnote{Various schools or authors attribute certain sets of characteristics to \Isvara. In our author's case, it is typically three: he often refers to \Isvara\ as all-knowing (\sarvajna), having all the power or capable of everything (\sarvasakti), and the \inidx{lord} of all (\sarvesvara) or the \inidx{supreme lord} (\paramesvara).  See for example, p.~\pageref{TTHSKT25} where the author constructs two \anumana\/s (inferences) that assert that \Isvara\ is \sarvasakti\ and \paramesvara, taking the one on \jnana\ in the \YBh\ as a model.}\\
  {}[He is] free from defects, [but] actions and rewards are superimposed on Him.\footnote{\label{OSTWO}The question whether \Isvara\ is capable of actions is discussed in conjunction with the question whether \Isvara\ has a body or not.  First the opponent asks the question in two paragraphs starting with \emph{tathāsarvajño} \reftoeditionpar{p25_58} (trl.\ p.~\pageref{TTHASRV25}), and the answer is given starting with the paragraph \emph{athāśarīratvād asarvajño} \reftoeditionpar{p25_130} translated on p.~\pageref{ATHSRRTVS}. There again the intended opponent is \Kumarila.

In this stanza the use of the word \emph{upahita}, deriving from the root \emph{upa-dhā-}, suggests the author's view that \kriya s and \phala s attributed to him are not real.  There is some ambiguity in the \phala\ part of this compound.  When the word appears in conjunction with the word \kriya\ (action), the most natural way to understand it is that it is the effect of an action.  Even though that may be figurative, it is hard to conceive that \Isvara\ receives a result of his action.  Another possible interpretation, adopted here, is that the \phala\ is the reward for actions of others.  The difficulty with this interpretation is that such concept is not mentioned in the body of the text.  Another problem with this compound, and in an effect, with the whole stanza is that the mention of defects, actions, and rewards with regard to \Isvara\ is redundant; it has already been mentioned in the previous stanza.  Note that this compound is not without textual problems.  The primary manuscripts, \Tm\ and L, both read \emph{°phalam} at the end of this line.  This reading is certainly not acceptable, and therefore a simple emendation, first introduced by the copyist of M, \emph{phalaḥ} is adopted here.  There remains some possibility that the reading \emph{phalam} was a result of further corruption.}%\\
\\  The Lord who is the cause of creation, end, and maintenance of everything,\\
  Reverence be to Him, the teacher even of the teacher.
\end{verse}\label{OSTZ2E}
%\section{First {Pāda}}

\section{Commentary on Pātañjalayogaśāstra 1.1}
\begin{quotation}
\st{Now\refm{4}%\pratika{atha}
\pratikap{p01_3} the instruction of yoga.~(\rYS{1.01}{1.1})}
\end{quotation}

\subsubsection{Explanation of the \sutra}

% \paragraph{On the word \emph{atha}}
The%\pratika{athetyādi}
\pratikap{p01_4}\labelh{ATHAATHATRL} \PYS\ starts with [the word] ``now (\atha{})'';  its explanation (\vivarana)\labelh{VIVPLNT} is commenced.\footnote{\label{STRMTSY}I have emended \emph{°śāstravivaraṇam} found in the manuscripts to \emph{°śāstram; tasya vi\-va\-raṇam}, which would be \emph{śāstrantasya vivaraṇam} without punctuation.  Considering that \emph{nta} and \emph{sya} look similar in the Malayalam script, this is a small emendation.  The adjective \emph{anā\-khyā\-ta\-samba\-ndha\-pra\-yo\-janam} in the next sentence refers to the noun \emph{°śāstram} (a neuter noun), rather than \emph{vivaraṇam}; this is indicated by the subsequent discussion on the relationship and the objective of a \emph{śāstra}.  The use of the adjective presupposes that the word modified by it has appeared earlier.  Note also that I read \emph{athetyādi} as a separate word, also modifying a neuter noun \emph{°śāstram}.

In addition, regardless whether this is \Sankara's composition or the one that imitates his style, many commentaries that are probably genuine \Sankara's, particularly those on the Upaniṣads, open up with the same formulation to the proposed reading here, viz., first a compound consisting of the first word(s) of the \emph{mūla} text and the word \emph{ādi}, followed by the title of the \emph{mūla} text, forming a phrase ``The TEXT TITLE, starting with the word(s) SUCH AND SUCH\ldots.''  That phrase is followed by the main clause, ``its commentary (\emph{vivaraṇa} or \emph{vivṛti}) is commenced (\emph{ārabhyate}).''  See the reference register in the edition (p.~\pageref{FIRSTSTC}).  That this beginning of the YVi is clearly intended to have the same basic structure (the first word of the \emph{mūla} text compounded with \emph{ādi}, followed by the title of the \emph{mūla}, and the phrase \emph{vivaraṇam ārabhyate}) is a support for my emendation.}  There (at the beginning of the explanation) if [the \sastra\ (discipline)] does not have the relationship (\sambandha) and the objective
(\prayojana) [explicitly] stated, it cannot [bring] a person to taking an action or backing off [from something].\footnote{The convention that a \sastra\
must have the objective (\prayojana), the relationship (\sambandha), and sometimes the \abhidheya\ is quite old.  Such discussions are found, for example, in \SVPRA{011--25}{kk.\ 11--25},
%\begin{quote}
%  \dnp{%
%    `अथातो धर्मजिज्ञासा' सूत्रमाद्यमिदं कृतम्। 
%    धर्माख्यं विषयं वक्तुं मीमांसायाः प्रयोजनम्॥११॥\\
%    सर्वस्यैव हि शास्त्रस्य कर्मणो वापि कस्यचित्। 
%    यावत् प्रयोजनं नोक्तं तावत् तत् केन गृह्यते॥१२॥ \\
%    मीमांसाख्या तु विद्येयं बहुविद्यान्तराश्रिता। 
%    न शुश्रूषयितुं शक्या प्रागनुक्त्वा प्रयोजनम्॥१३॥ \\
%    विद्यान्तरेषु नाप्येतद्यद्यभीष्टं प्रयोजनम्। 
%    अनर्थप्रापणं तावत् तेभ्यो नाशङ्क्यते क्वचित्॥१४॥ \\
%    मीमांसायां त्विहाज्ञाते दुर्ज्ञाते वाविवेकतः। 
%    न्यायमार्गे महान् दोष इति यत्नोपचर्यता॥१५॥ \\
%    तस्मात्प्रयोजनं पूर्वमुक्तं सूत्रकृता स्वयम्। 
%    यत् स्वेनोक्तं वदेयुस्तद्भाष्यकारादयः कथम्॥१६॥ \\
%    सिद्धार्थं ज्ञातसम्बन्धं श्रोतुं श्रोता प्रवर्तते। 
%    शास्त्रादौ तेन वक्तव्यः सम्बन्धः सप्रयोजनः॥१७॥ \\
%    शास्त्रं प्रयोजनं चैव सम्बन्धस्याश्रयावुभौ। 
%    तदुक्त्यन्तर्गतस्तस्माद्भिन्नो नोक्तः प्रयोजनात्॥१८॥ \\
%    सिद्धिः श्रोतृप्रवृत्तीनां सम्बन्धकथनाद्यतः। 
%    तस्मात्सर्वेषु शास्त्रेषु सम्बन्धः पूर्वमुच्यते॥१९॥ \\
%    यावत्प्रयोजनेनास्य सम्बन्धो नाभिधीयते। 
%    असम्बद्धप्रलापित्वाद्भवेत्तावदसंगतिः॥२०॥ \\
%    इह त्वाक्षिप्य सम्बन्धं भाष्य एवाभिधास्यते। 
%    धर्मप्रसिद्ध्यसिद्धिभ्यां तस्मान्नान्यो ऽभिधीयते॥२१॥ \\
%    न चाप्यत्राथशब्देन शास्त्रसम्बन्ध उच्यते। 
%    सम्बन्धक्रिययोर्ह्येष ब्रूते शास्त्राच्च ते पृथक्॥२२॥ \\
%    यो ऽप्ययं शास्त्रसम्बन्धो वर्ण्यते कैश्चिदादितः। 
%    क्रियान्तर्यरूपो वा गुरुपर्वक्रमो ऽपि वा॥२३॥ \\
%    तदतद्भावयोस्तस्य विशेषो नोपलभ्यते। 
%    श्रोतुर्विधौ निषेधे वा ज्ञाने वा शास्त्रगोचरे॥२४॥ \\
%    तस्माद्व्याख्याङ्गमिच्छद्भिः सहेतुः सप्रयोजनः। 
%    शास्त्रावतारसम्बन्धो वाच्यो नान्यस्तु निष्फलः॥२५॥  }
%\end{quote}
 \rnnNBh{1.1.1}/\rNV{1.1.1}.  The tradition may originally have been inspired by the opening discussion in the \VMBh\ although it discusses \emph{prayojana}, \abhidheya\ and \sambandha\ in a quite different context. \Sankara\ also held that a \sastra\ must have \sambandha\ and \prayojana.  For example, see the
introduction of the \rnnBhGBh{00intro}{intro.}, \rMaUBh{1.01}{1.1}, \rKaUBh{1.1.01}{1.1.1}, and \rMuUBh{1.1.01}{1.1.1}.  On the discussions of \emph{prayojana\/} in various Buddhist texts, see \citet{funayama:prayojana}.

\Sankara, like our author here, presupposes the view that the effect of a \sastra\ are \pravrtti\ (taking an action) and \nivrtti\ (backing off from something).  He, however, is critical of the view along with the related idea that a \sastra\ leads a person to \pravrtti\ and \nivrtti\ by means of \vidhi\ and \pratisedha.  See for example, \rBSBh{1.1.04}{1.1.4} \citep[particularly pp.~111\,ff., 130\,ff., 137\,ff.]{bsbh}.  There, \Sankara\ emphasizes that \sastra s do not have to have them as purpose; telling the truth without leading to any action can be the goal of \sastra s, especially the \Upanisad s.  In \rBSBh{1.1.01}{1.1.1} \citep[41--44]{bsbh} a somewhat related topic that \pramana\ leads to \pravrtti\ and \nivrtti\ is examined.  According to \Sankara, if the effect of \pramana\ is either \pravrtti\ or \nivrtti, human beings are no better than animals.%
% in the last reagrd, it might be helpful to take a close look at the Hetubindu and the Brahmasiddhi. Particularly the latter seems to offer a relevant discussion.
% % must have an objective and a relationship to human activities seems
% % to have been established as early as the time of NV. See NBh/NV
% % 1.1.1, {śV Pratijñāsūtra} kk.~11--25.  Also cf.\ the beginning
% % of the MBh, where \emph{proyojana}s are discussed.  On the
% % discussions of prayojana in various Buddhist texts, see
% % \cite{funayama:prayojana}.
}
Therefore the relationship and the
objective [of the \sastra], which are intended by the \Sutrakara, and which are the cause of commencement and termination of human [activities], are first discussed.

\paragraph{Objective}
Of\refm{6}%\pratika{tatra}
\pratikap{p01_5} those two, first, the \inidx{objective}---\labelh{PRYJN}

[It]%\pratika{cikitsāśāstravat}
\pratikap{p01_6} is explained [by the \Bhasyakara] by means of showing that it (the \sastra)
consists of \inidx{four divisions}, as in the \sastra\/ of healing (\cikitsasastra). %\footnote{On this discussion of
%  quadruple division see \cite{wezler:yvi2}.  Although Wezler ascribes
%  the beginning of the quadruple division to Buddhist, there are
%  evidences that suggest the existence of a text called \emph{%
%    Cikitsāśāstra}.  For example, a fragment (\emph{chāyā
%    madhuraśītalā}) quoted in YVi (p.~70, l.~5 of printed edition)
%  as from the \emph{Cikitsāśāstra} is also found elsewhere (in
%  \emph{Nyāyamañjarī}?  as a half verse and in \emph{%
%    Dvādaśāranayacakra}?  as a complete verse).  On the similar
%  discussion on \emph{heya}, etc., cf.\ NBh/NV 1.1.1.} 
Namely---the \sastra\/ of healing consists of four divisions: the disease,
the cause of the disease, freedom from disease, and the treatment; the goal [of this \sastra] is to explain the domains of the \inidx{four divisions} by means of a limited [set of] \inidx{injunctions} and \inidx{prohibitions}.\footnote{The phrase \emph{cikitsāśāstraṃ caturvyūham, rogo rogahetur ārogyam bhaiṣajyam iti} \index{caturvyuha@\emph{caturvyūha}} (including \emph{iti}) is found in \rYBh{2.15}{2.15}.  The editors of the 1952 edition suggested to emend the text by inserting \emph{tac ca} in between \emph{°dvāreṇa} and \emph{caturvyūhaviṣaya°}, making the sentence into two.  The text certainly reads more smoothly with the change, but it is still an intelligible sentence without the change; the text up to \emph{bhaiṣajyam iti} forms the subject and the rest predicate.  Even though the part \emph{cikitsāśāstram \ldots\ iti} is taken from the \bhasya, the author's intention was not to quote from the \bhasya, but use that phrase as part of his sentence.  The word \emph{iti} at the end of the phrase should be seen as terminating a list of four divisions.

The expression \emph{vidhi\-prati\-ṣe\-dha\-niyama\-dvāreṇa} is a result of an emendation proposed in the 1952 edition and adopted here.  It is a reasonable emendation, given a similar expression \emph{pra\-vṛtti\-ni\-vṛtti\-ni\-yama\-dvā\-re\-ṇa} later in the paragraph starting with \emph{yogānuśāsanam iti} \reftoeditionpar{p01_27} (translated on p.~\pageref{YGNSSM}).  It is a curious expression.  Literally the expression should mean ``by restricting\ldots.''}  In the same manner, here also, it is shown that the \sastra\ [of yoga] consists of four divisions [in the section] beginning with [the {sūtra}], ``For the one who knows the difference (\emph{vivekin}), everything is suffering, because of the sufferings from change, torment, and residues [of past experiences], and because of the conflict among the activities of the [three] \guna s (\rYS{2.15}{2.15}).''  That is:\refm{7} the object to-be-removed (\heya) is reincarnation (\samsara) full of suffering;\footnote{See \rYS{2.16}{2.16}: \emph{heyaṃ duḥkham anāgatam}.} its cause (\emph{hetu}) is the union of the seer and the to-be-seen (\emph{draṣṭṛdṛśyasaṃyoga}) that originates from \ignorance\ (\avidya);\footnote{See \rYS{2.17}{2.17}: \emph{draṣṭṛdṛśyayoḥ saṃyogo heyahetuḥ}, and \rYS{2.24}{2.24}: \emph{tasya hetur avidyā}.} the means of removal (\emph{hānopāya}) is the unwavering discriminative knowledge (\vivekakhyati);\footnote{See \rYS{2.26}{2.26}: \emph{viveka\-khyāti\-r a\-vi\-pla\-vā hāno\-pāyaḥ}.} and when the discriminative knowledge is present, \ignorance\ comes to an end. When it (ignorance) comes to an end, the union between the seer and the to-be-seen ends. This is the removal (\hana).  That verily is emancipation (\kaivalya).\footnote{See \rYS{2.25}{2.25}:
  \emph{tadabhāvāt saṃyogābhāvo hānaṃ tat dṛśeḥ
    kaivalyam}.  The expression in our text: \emph{duḥkha\-pra\-curaḥ saṃsāro heyaḥ; tasyā\-vidyā\-nimitto draṣṭṛ\-dṛśya\-saṃyogo hetuḥ; vi\-ve\-ka\-khyā\-ti\-r a\-vi\-pla\-vā hā\-no\-pā\-yaḥ; vi\-ve\-ka\-khyā\-tau ca sa\-tyām a\-vi\-dyā\-ni\-vṛ\-ttiḥ, ta\-nni\-vṛ\-ttā\-v ā\-tya\-nti\-ko dra\-ṣṭṛ\-dṛśya\-saṃ\-yogo\-pa\-ra\-mo hā\-nam} is a close parallel to that in \rYBh{2.15}{2.15}: \emph{tatra duḥ\-kha\-ba\-hu\-laḥ saṃ\-sā\-ro
    he\-yaḥ\dd\ pra\-dhā\-na\-pu\-ru\-ṣa\-yoḥ saṃ\-yo\-go he\-ya\-he\-tuḥ\dd\
    saṃ\-yo\-ga\-syā\-tya\-nti\-kī ni\-vṛ\-tti\-r hā\-na\-m\dd\ hā\-no\-pā\-yaḥ
    sa\-mya\-gda\-rśa\-nam\dd}.  Notable differences are the mention of \avidya\ as the cause of \samyoga\ and the replacement of \samyagdarsana\ with \vivekakhyati.  The first difference is a result of including the teaching of \rYS{2.24}{2.24} (\emph{tasya hetur avidyā}), and the latter incorporates the expression from \rYS{2.26}{2.26}.  The effect of these differences is that emphasis is put on the view that \avidya\ is the ultimate culprit.} As [this \sastra] aims at emancipation
which is comparable to freedom from disease [of the \sastra\ of healing], that very emancipation is its objective.

[Objection]\refm{8}%\pratika{nanu ca}
\pratikap{p01_7} The to-be-removed and its cause should not be
placed at the beginning, because [doing so is] pointless.  Because the
objective of the \sastra\/ is the removal.  [Thus] discriminative
knowledge, which is the means to that [aim (the removal)], only should be
discussed.  For, nobody enjoins the pain and its cause to someone
whose foot sole is pierced by a thorn, forgetting about the removal of
that [thorn].\footnote{Pain from a thorn is also found in \rNV{1.1.1} as an
  example of suffering (\emph{duḥkha}).  See the reference register accompanying this part of the text on p.~\pageref{NVKANTHA} for the text of the \NV\@.}

[Answer]\refm{9}%\pratika{naitad}
\pratikap{p01_8} That is not correct.  Because the means of removal
presupposes the to-be-removed and its cause.  It is not possible to
discuss discriminative knowledge, which is the opposite of 
\ignorance, until it is stated that this cycle of reincarnation is what 
should be removed by this means of removal, and that its cause is the
union---itself caused by \ignorance---between the seer and the
to-be-seen.  Because we also\footnote{The presence of the word \emph{api} here is conjectural. Seeing various corrections both in \Tm\ and L, I infer that there was a confusing correction already in the archetypal manuscript. The exact situation in the archetype is hard to reconstruct but I consider the presence of \emph{vi} in \Tm\ significant. In the Malayalam script \emph{p} and \emph{v} are often hard to distinguish. The presence of the word \emph{api} fits in the context. The subject matter of this phrasal clause expressing a reason is not identical to the preceding. The reason rather functions by pointing out a real word example. In that sense, qualifying the subject with ``also'' is appropriate.} experience that a patient, who has something
to-be-removed and its cause, and who seeks for the means of removal
[of the disease] goes for treatment.  For, the \sastra\/ of healing is not taught without presupposing disease and the cause of disease.

\paragraph{Relationship}
The\labelh{SMBDHPI}%\pratika{sambandho 'pi}
\pratikap{p01_9} relationship, too [is properly presented in the \sastra]---

The%\pratika{vivekakhyāter}
\pratikap{p01_10} result (\emph{phala}) to be achieved (\emph{sādhya}) by discriminative knowledge is but the removal; also, the means (\emph{sādhana}) of removal is but discriminative knowledge.  Precisely this mutual determination (\emph{itaretaraniyama}) of the aim (\emph{sādhya}) and the means (\emph{sādhana}) [is the relationship]; there is no other relationship.  This is analogous to the mutual determination that the result of the treatment is but the freedom of disease and that the means toward freedom of disease is but the treatment.  %Similarly, the result of the treatment is but the freedom of disease, and the means toward freedom of disease is but the treatment.  Therefore they restrict each other. 
 And this [mutual determination] is based on the \sastra.\footnote{The word \emph{iti} at the end of this paragraph marks the end of the discussion on \prayojana\ and \sambandha\ of this \sastra.  The use of \emph{iti}, followed by a concluding remark that starts with the word \emph{tasmāt} is found frequently in this text.  See another example on p.~\pageref{TSMDTH}.}

In%\pratika{tasmāt}
\pratikap{p01_11} conclusion,\refm{11} the instruction of yoga (\yoganusasana) [that is this \sastra] has proper relationship and objective.\labelh{PSCOMP}

\paragraph{Justifying the first \sutra}\labelh{JSTF}
[Objection]\refm{12}%\pratika{nanu ca}
\pratikap{p01_12}\footnote{Our author starts discussions to justify why the reading \emph{atha yogānuśāsanam} is the first \sutra, raising several other alternatives for the \sutra.  I doubt that there were real criticisms on the wording of the first \sutra\ of the \PYS.  The following arguments rather appear to be purely an intellectual exercise to write a commentary that guards the text being commented.} If the removal is the
objective, and its means is discriminative knowledge, then [the first \sutra] should say ``Now the instruction
of discriminative knowledge (\emph{atha
  vivekakhyātyanuśāsanam}).''  For what purpose does the [first]
\sutra\ say ``Now the instruction of yoga (\emph{atha
  yogānuśāsanam})''?

[Answer]%\pratika{tadupāyatvād}
\pratikap{p01_13} Because\refm{13} yoga is the means of
obtaining it (\vivekakhyati).  Precisely the means should be
stated.  For, if the goal (to be obtained, \emph{upeya}) is [first] explained, then since the
goal presupposes the means (to obtain it), it will furthermore be
necessary to speak of the means along with its divisions and
in its entirety.  On the contrary, when it (the means, yoga) is denoted,
then everything is denoted.

[Objection]\refm{14}%\pratika{kathan tad°}
\pratikap{p01_14} How is [yoga] the means to obtain that (\emph{vivekakhyāti})?

[Answer]%\pratika{yata aaha}
\pratikap{p01_15} Because\refm{15} the \Sutrakara\ says: ``When
impurities are removed as the result of exercising the divisions of
yoga, the intelligence begins to shine until [it becomes]
discriminative knowledge (\rYS{2.28}{2.28})''; ``When skilled in \emph{%
  nirvicāra[samādhi]}, there will be internal tranquility (\rYS{1.47}{1.47})'';
and ``Then there will be an intellect called carrier of the truth (\emph{%
  ṛtambharā}) (\rYS{1.48}{1.48}).''  Also the \Bhasyakara\ says:
``That (\emph{samādhi}), % needs to be reconsidered
when the mind is focused, illuminates the real object \ldots\ (\rYBh{1.01}{1.1}).''  And discriminative knowledge is the understanding
of the real object.  Accordingly, it is appropriate that the instruction of yoga---the means to obtain it (\emph{vivekakhyāti})---has been put in a sūtra at the beginning.

[Objection]\refm{16}%\pratika{nanu ca}
\pratikap{p01_16} Is it not that discriminative
knowledge arises after exercising the divisions of yoga? Then [the first
\sutra] should state ``Now the instruction of the divisions of
yoga (\emph{atha yogāṅgānuśāsanam}).''\footnote{The \PYS\ in fact teaches the eight divisions of yoga (\emph{yogāṅga}s), starting from 2.28.}

[Answer] No.\refm{16+1}%\pratika{na}
\pratikap{p01_17} Because [a \sastra\/
should have] a beginning with the result.\footnote{The answer in this paragraph, which may be paraphrased as ``it is appropriate to start a \sastra\ by stating the result'' may seem contradictory to what has just been argued: the means to obtain a goal, rather than the goal should be taught.  As can be observed in the paragraphs that follow, our author distinguishes a goal (\emph{upeya}) and the result (\emph{phala}).  That seems to be his strategy to avoid such a contradiction.  Still, the arguments appear weak.  This is one reason to suspect that this part was an intellectual exercise to guard the text being commented.}  For, since yoga is the result of
exercising divisions of yoga, it is appropriate to begin with that.

[Objection] If%\pratika{yady evaṃ}
\pratikap{p01_18} so,\refm{17} the beginning should be the removal.

[Answer] No.%\pratika{na}
\pratikap{p01_19}\refm{17+1} Because that is a goal.  Namely, it is nothing but
a goal.  On the other hand, yoga is the goal (\emph{upeya}) for its own divisions, and it is also the means to obtain (\emph{upāya}) discriminative knowledge and [supernatural powers, such as] the ability to become subtle (\emph{aṇimā}),\footnote{In the \PYS\ \rYS{3.44}{3.44} and accompanying \bhasya\ mention the well-known supernatural powers \emph{aṇimā}, \emph{mahimā}, etc.} and so on.  Therefore it is more valuable, because it is the cause of two fruits.\footnote{The reading \emph{phaladvayanimittatvāt} is conjectural.  The readings in two primary manu\-scripts (\Tm\ and L) are unintelligible.  It is not even certain what the scribe of L meant to write because there the letters \emph{pa} and \emph{va} are often indistinguishable.  Of the three possible readings found in the primary manuscripts (\emph{phalavayanimitta°} or \emph{phapapayanimitta°} for L and \emph{phalavannimitta°} for \Tm), \emph{phalavayanimitta°} requires the least change to produce a reasonable reading, \emph{phaladvayanimitta°}, which I have adopted.  We can be reasonably certain that the corrupt reading \emph{phalavayanimitta°} was already present in the common ancestor (possibly the direct exemplar of both \Tm\ and L).  The reading found in \Tm\ was probably a result of the scribe's attempt to produce a meaningful reading.}

[It is\pratikap{p01_20} not\refm{18} appropriate to have the
\sutra, ``Now the instruction of the divisions of yoga,'' at the
beginning]%\pratika{uttarasūtrārtha°}
 also because [the first \sutra] introduces the next \sutra.
For, if the first \sutra\ was to say ``Now the instruction of the
divisions of yoga,'' then it would not be appropriate to say
``[It (\emph{asamprajñāta-samādhi}) is] yoga [and it is] the termination of the activities of the mind (\rYS{1.02}{1.2}).''\footnote{\label{YS12first}The translation here follows the interpretation of the \sutra\ in the YVi where the author states that the \sutra\ is the definition of the \emph{asamprajñāta-samādhi}.  This is somewhat counter-intuitive, but our author is clear about \sutra\ 1.2 being the definition of \emph{asamprajñāta-samādhi}.  See, for example, the last two paragraphs of the commentary of the current \sutra.  See also note \ref{no_yogah}.}
In that case, the [next] \sutra\ should read ``\emph{yama, niyama},''
and so on (\rYS{2.29}{2.29}).

If%\pratika{atha tad eva}
\pratikap{p01_21} you say---[Objection]\refm{19} But then that is exactly how [the \sutra] should read---[Answer] No.  For the next \sutra, ``[It (\emph{asamprajñāta-samādhi}) is] yoga [and it is] the termination of the activities of the mind,'' has the purpose of showing the relationship
between discriminative knowledge---whose means is yoga,\footnote{The editors of the 1952 edition suggested the reading \emph{yogopeyāyā(ḥ)} instead of \emph{yogopāyāyā(ḥ)}, the reading attested in all the manuscripts.  It is understandable that they preferred a \TP\ compound rather than a \BV\ compound, but the following word, another adjective to modify the word \emph{vivekakhyāteḥ}, reads \emph{hānopāyābhūtāyā(ḥ)}.  The author made an extra effort to makes it clear that \emph{hānopāya} is to be read as a \TP\ compound.  I take it that the author was illustrating the position of \emph{vivekakhyāti} using two compounds that both have the word \emph{upāya} as a member.  Even without the use of \emph{°bhūta}, in the second compound the case ending of the word \emph{yogopāyāyāḥ} makes it clear that it is a \BV\ compound.}
and which is the means to obtain the removal---and the removal which is
the result.

[Objection]\refm{20}%\pratika{katham punar}
\pratikap{p01_22} How, then, can it
have the purpose of showing the relationship while it states the
definition (\emph{lakṣaṇa}) of the seedless immersion (\emph{%
  nirbīja\-samādhi})?\footnote{Our author considers \rYS{1.02}{1.2} to be the definition of \emph{nirbīja-samādhi}/\emph{asamprajñāta-sa\-mā\-dhi}.  See note~\ref{no_yogah}. Cf.\ also note \ref{YS12first}.}
  % SUPPLY THE REFERENCE!

[Answer]\refm{21}%\pratika{satyam evam}
\pratikap{p01_23} That is correct.  Even so, the \sutra, which
is [in fact] only setting forth the relationship between the removal and its means,
amounts to the definition of seedless immersion (\emph{%
  nirbīja\-samādhi}).  For, the removal is not different from the
immersion which is characterized as suppression  (\nirodha).
There still is the following distinction: in the case of immersion
characterized as suppression, there will be further
commencement (\emph{pravṛtti}); on the other hand, in the case of removal, cessation (\emph{nivṛtti}) is permanent.  But in the state of immersion [characterized as suppression] there is no distinction from the removal at all.\footnote{This paragraph offers a glimpse of our author's soteriological position, which is more clearly stated in two paragraphs below. While defending the arrangement of the \PYS, he implies that the best yoga can offer, \nirbijasamadhi, is still inferior to removal (\hana).  Discriminative knowledge is, as repeatedly expressed, is the means to achieve the final goal.  What one might consider as the highest goal in the system of the \PYS, \nirbijasamadhi, is just a temporal experience of real salvation, according to our author.}

Namely,\refm{22}%\pratika{tathā cāha}
\pratikap{p01_24} [the \Sutrakara] says [in the immediately following \sutra\ after 1.2], ``At that time the seer stays in his original state (\rYS{1.03}{1.3}),''  [using the same expression as in the last \sutra] ``Or\footnote{I have emended the particle \emph{ca} found in all the manuscripts and in the 1952 edition to \emph{vā}, following the reading of the \sutra.  There is no indication in the YVi manuscripts that our author had the reading \emph{ca} for \sutra\ 4.34 when it is quoted and explained.  The exchange may have been introduced by a scribe who was unaware that the particle \emph{vā} was part of the quotation but thought that there should be \emph{ca} to connect two sentences.

The editors of the 1952 edition proposes to insert \emph{iti} after the first quote (from \rYS{1.03}{1.3}).  I do not find it necessary.} the power to cognize (\citisakti)\footnote{In the \PYS, \purusa\ is identified as \citisakti.  In the YS the word \citisakti\ appears only in the last \sutra\ 4.34, but it already appears in the \bhasya\ on \sutra\ 1.2.  The author of the \bhasya\ continues to paraphrase \drastr\ in \rYS{1.03}{1.3} with \citisakti.  The author of the YVi glosses the word as a \KD\ compound (\emph{citir eva śaktiś citiśaktiḥ}) \citep[11,26]{yvi}.  Also, he understands the part \emph{citi} denotes the verbal root \emph{cit-} and hence the \sakti\ to be the one expressed by the verbal root.  This observation is based on the sentence \emph{yathā pacyādayaḥ śaktayaḥ śaktimadapekṣāprāpaṇīyajanmāno na svayaṃ śaktayaḥ} that immediately follows the compound analysis above.  That \purusa\ is the root \emph{cit-} (following the understanding of the author of the YVi) is repeatedly expressed in the \YBh\ and in the YVi (first expressed in the \rYBh{1.02}{1.2} \emph{citiśaktir apariṇāminī\ldots} on which the YVi's compound analysis above is given).  The view implicit in such an interpretation is that \purusa\ is changeless because it is a verbal root \cit- and because the capacity (\sakti) of a verbal root is consistent.  Similarly \rYS{2.20}{2.20} has \emph{draṣṭā dṛśimātraḥ}, again identifying \purusa\ with a verbal root \emph{dṛś-}.  There it is more explicit that the root \emph{dṛś-} exists without a substance or a subject.  This is a solution to the problem that anything denoted by a substantive is subject to change.  This makes us think of \inidx{grammarian}s' influence on the formation of the \PYS.  It should be reminded that the beginning of the \PYS\ imitates that of the \VMBh\ and that the \PYS\ accepts the \sphota\ theory.  Note also that identification of \purusa\ (masculine) with \citisakti\ (feminine) is unique to the \PYS, absent in the \Samkhya\ tradition.  Another interesting point is that both in the YBh and the YVi the root \cit- is seen as special in that it does not require the locus. This is expressed in the YVi on \rYS{1.02}{1.2} \emph{citiḥ punaḥ svayam eva śaktir nārthāntaram apekṣate}, following the sentence \ldots\ \emph{svayaṃ śaktayaḥ} above.  The view that \purusa\ is nothing but the root \emph{cit-} is presupposed in the discussion in \rYBh{1.03}{1.3} starting with \emph{yadi citir eva puruṣaḥ\ldots}.  Such discussions are echoed in \Sankara's discussion on the meaning of the root \emph{upalabh-} in the \Upad\ 2.2.76i\,ff.\ where \Sankara\ argues that \atman/\emph{upalabdhṛ} is eternally only \emph{upalabdhi} (\emph{nityopalabdhimātra eva hy upalabdhā}), and that \atman\ does not change.} being established in its original state (\rYS{4.34}{4.34})'' is emancipation.\footnote{Although the last word of this sentence (\emph{kaivalyam}) is not enclosed in the quotes, it is necessary to complete the sentence meant to be referred. The whole \sutra\ 4.34 reads: \emph{puruṣārthaśūnyānāṃ guṇānāṃ pratiprasavaḥ kaivalyaṃ svarūpapratiṣṭhā vā citiśaktir iti}.  The word \emph{kaivalyam} is implied to be supplied in the second statement.  It is important for our author to have the word \kaivalya\ in this sentence to illustrate that similar expressions are used in the YS to describe the state in \asamprajnata\samadhi\ and \kaivalya.}
For, emancipation, [characterized as \citisakti's] being established in its original state, is the removal (\hana).\footnote{Two emendations are involved in this sentence.  The first is the reading \emph{sva\-rūpa\-prati\-ṣṭha\-tvaṃ} adopted from the 1952 edition instead of \emph{°prati\-ṣṭhaṃ} in the manuscripts.  The former reading is preferable because what we want here is a substantive rather than an adjective.  The other emendation is from \emph{kaivalyam āha} to \emph{kaivalyaṃ hānam}.  Considering that this paragraph elaborates the point at the end of the last paragraph (``in the state of immersion (\samadhi) there is no distinction to the ultimate removal (\hana),'') there should be a reference to \hana.  The reading, found in all the witnesses, is non-sensical, but the presence of the syllable \emph{ha} in that reading suggests that it is a corruption involving the word \hana.  Our author is pointing out the state in \asamprajnata\samadhi\ (defined in \rYS{1.02}{1.2}) is described in the same way as \kaivalya\ is described in the YS and therefore that they are not different.}  Therefore, in order to strengthen the
understanding of the purpose of the \sastra, the emancipation is revealed by [defining] the seedless immersion (\emph{nirbījasamādhi}) [in the next \sutra].

Accordingly,\refm{22+1}%\pratika{tasmād yad}
\pratikap{p01_25} it is not true that seedless immersion
(\emph{nirbīja\-samādhi}) is the means to establish
emancipation, as some say.\footnote{Cf.\ \Vacaspati\ (\Tattvavaisaradi) on \rYBh{1.01}{1.1} \emph{%
    niḥśreyasasya hetuḥ samādhir iti hi
    śruti\-smṛtī\-tihāsa\-purāṇe\-ṣu pra\-siddham\dd} \citep[2,4--5]{ys:anss}.  However, such a view that \nirbijasamadhi\ brings about the emancipation is a straight-forward interpretation of the beginning of the \PYS.  Our author's interpretation that (discriminative) knowledge is the means to attain emancipation is rather forced and particular.  This might be comparable to similar insistence on knowledge in \Sankara's works, such as the \BSBh, \BhGBh, etc.}  Instead, only knowledge is the means, through
cessation of \ignorance.  Because the state of being bound [to the cycle of rebirth] is caused by ignorance.
Therefore even if ``the instruction of yoga'' is placed in the [first]
\sutra, since the purpose of yoga is knowledge, [the \Sutrakara]
shows the relationship between knowledge and removal, too, with this
[\sutra].  Therefore [the \sutra\ is] appropriately composed.  Also, in order [to express] the relationship with the next \sutra, the [first] \sutra\ should contain [the words] ``the instruction of yoga.''\footnote{This paragraph concludes the discussion on the appropriateness of the first \sutra\ started in the paragraph \emph{nanu ca} \reftoeditionpar{p01_12} translated on p.~\pageref{JSTF}.  The argument in this paragraph is a little forced.  Still, one may note the author's insistence that knowledge is the means.}

\medskip
Even\refm{24}%\pratika{yady api}
\pratikap{p01_26} though there is also a teaching of the
means for those who seek a result (\emph{phala}) and have come
to expect the means to obtain the result, this (the means) is not the
objective.   Since everyone acts in view of a result, that (the
result) is the sole objective.\footnote{I have not given a literal translation of the word \emph{iti} at the end of this sentence.  The word \emph{iti} and even the whole sentence appear to be a little out of place.  Most likely the word \emph{iti} here signals the completion of a major segment of the text.  And the segment could be the preliminary discussion on the objective and the relationship of the whole \sastra.  However, the discussion on the objective and the relationship had completed (see page \pageref{PSCOMP}), and the mention of only \prayojana\ here is is not congruent with the immediately preceding discussions.  It is possible that this is a trace of a correction going wrong; this passage was originally omitted while copying but supplied in the margin with a marking in the body indicating where to insert the text; the insertion point was misunderstood by the next scribe. If this passage is indeed misplaced, the proper place would be at the end of the discussion on \prayojana, i.e., immediately before the passage starting with \emph{sambandho 'pi} \reftoeditionpar{p01_9} translated on p.~\pageref{SMBDHPI}. The presence of \emph{iti} at the end of this passage also becomes explicable if it indeed is misplaced. The word \emph{iti} would mark the end of the discussion on \prayojana.}

\subsubsection{Explanation of the \bhasya}
\label{YGNSSM} ``The%\pratika{yogānuśāsanam iti}
\pratikap{p01_27} instruction of yoga''---\refm{25} This is comparable to [the way] a student is instructed by means of 
determining specific actions (\emph{pravṛtti}) [that he should do] and withdrawal from actions (\emph{nivṛtti}) [that he should not do]. Likewise [the \sastra] is called the instruction of yoga, like the instruction to pupils  (\emph{antevāsyanuśāsana}), since it essentially resembles [the instruction to pupils] in that [the \sastra] determines a specific goal (\emph{sādhya}), means (\emph{sādhana}), and their divisions (\emph{aṅga}).\footnote{I suspect that the sentence, especially the part \emph{tathā viśiṣṭa\-sādhya\-sādhana\-tad\-aṃga\-niyama\-mātra\-sādṛśyāt}, is corrupt.  Given the somewhat similar discussion in the \TaiUBh\ (see the apparatus on p.\ \pageref{TAIUBH111REF} for the text), the discussion could have been longer.  However, I cannot offer a viable emendation.%  Note the expression \emph{°dvāreṇa} and the use of the word \emph{niyama} in the \TaiUBh.  Our author appears to be fond of these expressions in this part of the commentary.  They disappear in the rest of the work)
}  ``Instruction'' = instructing (\emph{anuśiṣṭi}).\footnote{This seemingly meaningless—particularly in translation—paraphrase is meant to reject the ambiguity that the word \emph{anuśāsana\/} with the \emph{-na\/} suffix can denote the means to instruct.  By paraphrasing \emph{anuśāsana\/} with \emph{anuśiṣṭi}, the author claims that the word \emph{anuśāsana\/} is meant to refer to the action of instruction, not the means of instruction.}
\bht{The \sastra\ is the instruction of yoga}\footnote{\label{VMBHEMU}I consider the sentence or the phrase \emph{yogānuśāsanaṃ śāstram} to be part of the \bhasya.  This is equivalent to the sentence \emph{yogānuśāsanaṃ śāstram adhikṛtaṃ veditavyam} in most transmission of the \YBh\ after the sentence \emph{athety ayam adhikārārthaḥ} (see below).  The sequence, starting from the \sutra, \emph{atha yogānuśāsanam\ddd\ athety ayam adhi\-kārā\-rthaḥ\dd\ yogānuśāsanam śāstram adhikṛtaṃ veditavyam} forms a closer parallel to the opening of the \VMBh: \emph{atha śa\-bdā\-nu\-śā\-sanam\dd\ athety ayaṃ śabdo ’dhikārārthaḥ prayujyate\dd\  śa\-bdā\-nu\-śā\-sa\-naṃ śā\-stram adhi\-kṛtaṃ ve\-di\-tavyam}.  In my reconstruction the \sutra\ and \bhasya, combined, reads \emph{atha yo\-gā\-nu\-śā\-sa\-nam\ddd\ yo\-gā\-nu\-śā\-sanam śā\-stram\dd\ athety ayam adhi\-kārā\-rthaḥ\dd}.  It is conceivable that the version of the \YBh\ our author used had the opening not exactly parallel to that of the \VMBh. It may, however, be noted that there are some uncertainties how well the text of the \VMBh\ is established in the editions. Cf.\ \citet[248--50]{cardona:1976} for the overview of the discussions on the transmission of the \VMBh. Cf.\ also \citet{witzel:1986}. Note that the reading of the quote from the \VMBh\ in the YVi below (see pp.\ \pageref{VMBH1138_E} for the edition and \pageref{VMBH1138_T} for the translation) is somewhat different, although close enough to call it a quote, from the readings found in editions of the \VMBh. Still, the parallelism is easily recognizable.  The current paragraph in our text could be easily seen to be an explanation of the word \emph{yogānuśāsanam} in the \sutra\ itself.  However, it is preferable to see that our author has started to gloss the \bhasya, and its first sentence is \emph{yogānuśāsanaṃ śāstram}.  Otherwise, the presence of the word \emph{śāstram} is inexplicable.  The sentence should be seen as \bhasya's gloss of the word \emph{yogānuśāsanam} in the \sutra.} by which or in which yoga is instructed.%\footnote{Even though it is not clear from the translation, the syntax in the YVi suggests that the author considers that in the paraphrase in the YBh (\emph{yogānuśāsanaṃ śāstram}) the word \emph{yogānuśāsanam\/} modifies the word \emph{śāstram\/} as a \emph{bahuvrīhi\/} compound.  According to this interpretation, the \Bhasyakara\ interprets the beginning of the YS as proclaiming the topic by introducing the word \sastra\ as implied in the \sutra.}  %[The compound] \emph{yogānuśāsana} means either that yoga is
%instructed by it, or in it, and it is the \sastra.

\bht{That%\pratika{athety ayam}
\pratikap{p01_28} word\refm{27} ``now'' means
  the declaration of the topic.}  %\footnote{Note that the first \sutra\ and this first
  %sentences of YBh is parallel to the beginning of MBh:
%
%\bigskip
%\begin{tabular}{p{.5\textwidth}p{.5\textwidth}}
%  \multicolumn{1}{c}{YS + YBh} & \multicolumn{1}{c}{MBh}  \\
%  \emph{atha yogā\-nu\-śā\-sa\-nam| athe\-tya\-ya\-m
%    a\-dhi\-kā\-rā\-rthaḥ| (yogā\-nuśā\-sa\-nam adhi\-kṛtaṃ
%    ve\-di\-ta\-vyam||)} This last sentence was not present in the YBh
%  commented by the author of the YVi.  See note\fnref{note:veditavyam} on
%  page~\pageref{note:veditavyam}.  & \emph{atha
%    śabdā\-nu\-śā\-sa\-nam| athe\-ty ayaṃ śa\-bdaḥ
%    adhi\-kā\-rā\-rthaḥ pra\-yu\-jya\-te|
%    śa\-bdā\-nu\-śā\-sa\-naṃ śā\-stra\-m adhi\-kṛ\-taṃ
%    ve\-di\-ta\-vyam|}    \\
%\end{tabular}
%\bigskip } 
``The declaration of the topic (\emph{adhikāra})'' = opening (\emph{%
ārambha}) = announcement (\emph{prastāva}), and ``meaning (\emph{%
artha})'' = sense (\abhidheya).  [These words are connected as
a \KD\ compound and then form a \BV\ compound.]\footnote{The sentence is a typical \emph{vigraha\/} to analyze a compound, \emph{adhikārārthaḥ} in this case.  The author of the YVi is clarifying the meanings of constituent words \emph{adhikāra\/} and \emph{artha} by listing synonyms because they both have divergent meanings.  Also, he is showing the relationship between the constituents (in the same case, viz., \KD) and the type of the compound syntactically.}  [This interpretation is]
based on the authority of a \smrti\/ (the Mahābhāṣya) by hands of a learned.\footnote{\label{PTJME}Thus the learned person is \Patanjali.  \Patanjali, the author of the \VMBh, is mentioned by name in the YVi in the commentary on \rYS{3.44}{3.44} \citep[299]{yvi}.  However, our author shows no awareness of \Patanjali\ being associated with the text he comments upon, viz., the \PYS.  Note that the beginning of the \PYS\ emulates that of the \VMBh.  See the parallel passage in the edition part.  See also note \ref{VMBHEMU}.}%For the relationship between the \emph{śiṣṭa\/}s and \emph{śāstra}, see the opening discussion of the Mahābhāṣya itself.  There Patañjali declares that ...}
  
%\footnote{Is he referring to the MBh?}

[Objection]\refm{28}%\pratika{nanu cā°}
\pratikap{p01_29}\footnote{Our author consistently starts the objection with the expression \nanuca\ as found here.  As noted by \citet[36]{wezler:yvi1}\index{Wezler, Albrecht}, this is fairly consistent in the YVi and might indeed be idiosyncratic\index{idiosyncrasy} to the author.  A notable text that uses the expression \nanuca\ frequently is the \VMBh\ that set the standard for Sanskrit commentaries.  In \Sankara's BSBh the use is significantly less (as far as we look at published editions).  Still, the Nirnaya Sagar Press edition of the BSBh has the expression in 2.3.19, 2.3.43, 2.3.47, 4.2.5, 4.2.6, and 4.2.20.  The concentration is a little intriguing.  It is nonetheless necessary to wait for a critical study of the \BSBh\ manuscripts before concluding anything from the use of this expression.  We cannot exclude the possibility that the \BSBh, being one of the most heavily and closely studied text, has gone through many generations of cleanup in style.} The learned ones teach in a \smrti\/ that the word \atha\ means ``immediately after.''  Namely,
[\Sabara] says:\label{SBRATH}
\begin{quote}
  It is experienced that [this word \emph{atha\/} at the
  beginning of the Mī\-māṃ\-sā\-sū\-tra (\JS)\footnote{\rJS{1.1.1}: \emph{athāto dharmajijñāsā}.}] has the meaning of introducing something
  that comes immediately after something has concluded.\footnote{\rSBh{1.1.1}.  The
    whole sentence runs as follows: \emph{tatra loke 'yam athaśabdo
      vṛttād a\-na\-nta\-ra\-sya prakriyārtho dṛṣṭaḥ\dd.} A similar discussion
    of the meaning of \atha{} at the beginning of \sutra\ texts
    appears also in the \rnnBSBh{1.1.01}{1.1.1} on \rBS{1.1.01}{1.1.1}: \emph{athāto brahmajijñāsā} (note the parallel structure of the \rJS{1.1.1} and the \rBS{1.1.01}{1.1.1}).  \Sankara\ holds that the meaning of \atha{} at the beginning of BS 1.1.1 is \emph{ānantarya}, just like Śabara, although he claims the intention to learn about Brahman (\emph{brahmajijñāsā}) does not necessarily come after the intention to learn about Dharma (\emph{dharmajijñāsā}).  The antagonist of \Sankara\ here holds the view that the \brahmajijnasa\ must come after \dharmajijnasa.  \Sankara\ was probably dealing with the traditional view that the two \sutra\ texts (the JS and the BS) formed a set and to be studied in succession.  (\Sankara\ considers the \JS\ and the \BS\ form one continuous text; he refers to \rJS{3.3.14}: \emph{śruti\-liṅga\-vākya\-prakaraṇa\-sthāna\-samākhyā\-nāṃ sama\-vā\-ye pāra\-daurbalyam artha\-vi\-pra\-karṣāt} as belonging to ``an earlier section (\emph{pūrvasmin kāṇḍe}).'') As \emph{dharmajijñāsā} is placed at the beginning of the JS, this interpretation of the meaning of \atha{} allows one to begin studying the \BS\ without learning the JS, or the discipline of Mīmāṃsā.  We may say that \Sankara\ and the author of the YVi are in agreement in that they both refer to the \Mimamsaka\ view of the interpretation of the word \atha\ and that they are both anti-\Mimamsaka\ when the word \atha\ is concerned.  I am unaware of any other author who is so concerned about how the \Mimamsaka s understand the meaning of the word \atha\ as \Sankara\ or the author of the YVi.  Note that the author of the YVi here ignores the more developed discussion by \Kumarila\ in his \SV\ while \Sankara\ appears to be aware of it.}
\end{quote}

[Answer]\refm{29} No.%\pratika{na\dd}
\pratikap{p01_30}  Because [the word
\atha{}] always has the meaning ``opening,'' and because the
meaning ``immediately after'' must be [secondarily] understood.  For
example---when we say ``[he is] a son,'' it is understood that [the man can be] a father.\footnote{The sentence is vague, but the idea behind this example should be the one that frequently appears in śāstric literature---that a man can be a son or a father or many other things.  In that vein, the text would read better if we replace \emph{putra} with \emph{puruṣa} and \emph{pitṛ} with \emph{putra}. (A man is always a son.) There is some suspicion that the original text read as such and the reading we have is a result of corruption. The idea is again already expressed in \rVMBh{1.1.66}{1.1.66}.  \Sankara\ mentions the idea in \rBSBh{2.2.17}{2.2.17}.  (See page \pageref{PTRPITR} for the text.)  There, \Sankara\ refers to the written symbol of the number one in the decimal notation system.  He says it is a line (\emph{rekhā}).  The passage is a paraphrase of another in \rYBh{3.13}{3.13}.  Here we might have a testimony to the writing system (a line representing the numeral 1) \Sankara\ and the author of the YBh were familiar with.}  But the meaning of ``son'' is not ``father.''  In the same
manner, here, only ``opening'' is being directly expressed.  On the other
hand, the meaning ``immediately after'' is [secondarily] understood.
For this very reason, [Śabara] says ``It is [often] experienced that
[this word \atha{} at the beginning of the JS] has the meaning of
introducing something that comes immediately after something has concluded.''  If the
meaning of \atha{} was``immediately after,'' then he would have
said [more explicitly], ``[This word \atha{} is used] to signify
that the beginning comes immediately after something has concluded.''
%[instead of saying ``it is experienced \ldots,'' in which case it is
%implied that that is not always the case.

Also,%\pratika{anantarasyeti}
\pratikap{p01_31} [the word] ``something that comes immediately after (\emph{anantara})''  signifies the thing[, a substance,] that comes immediately after (\emph{anantarabhāvin}).  %[This is the evidence that even \Sabara\ did not consider that the meaning of the word \atha\ was `immediately after' since \Sabara\ does not use the expression `immediately after (\emph{ānantaryārthe}),' but uses `something that comes immediately after (\emph{anantarasya}).']
% When I was trying to locate the usage of the simil putra-putṛ, I 
% encountered a reference to Vaiśeṣika-s in \rBSBh{2.2.17}{2.2.17}, where 
% śa"nkara lists six padārthas of V.  Was there V at the time of 
% S?  Also there is a mention to satkāryavādin that include Manu.  
% What is he referring?  Off course there is a mention to Sāṃkhya.  
% Who are the yogins for S?  Who are the Vedavādins?  Remember 
% Vedavādins in Jayamangalā? [This turns out to be Māṭharavṛtti 20080827]
% 
And,\refm{30} %%\pratika{api caivaṃ}
% \pratikap{p01_32}
a \smrti\/ (\VMBh) indeed says thus:
\begin{quote}
  The primary meaning of some indeclinables is nominal, of others it
  is primarily verbal.  Such words as ``upwards (\emph{uccair}),''
  ``downwards (\emph{nīcair}),'' etc., have the nominal meaning as
  their primary meaning.  Such words as ``near (\emph{hiruk}),''
  ``separately (\emph{pṛthak}),'' etc., have the verbal meaning as their
  primary meaning.  There is no other meaning of indeclinables apart from these. (\rVMBh{1.1.38 or 1.2.47}{1.1.38 or 1.2.47})\footnote{Note that the quote here is not exactly the same as what we find in edition of the \VMBh. Parallel passages appear twice in the \VMBh, and in both cases, the editions of the \VMBh\ have \emph{ādi} at the end of examples while the quote here does not. Perhaps more importantly, the whole part on \emph{taddhita} or on \emph{tibanta} is missing in the quote here in the YVi. Instead, it has a sentence to exclude the third category of \avyaya s.}\labelh{VMBH1138_T}
\end{quote}
Hence, if the meaning of the word \atha\ is ``immediately after'' [as you
claim], hearing a nominal suffix is reasonable, because its primary
meaning is a substance.  On the other hand, if [the meaning of \atha\ is] the action of beginning [as I claim], then it is reasonable not to hear
a nominal suffix, because its primary meaning is something that does not
exist[; and that is the case].\footnote{As seen in the critical apparatus, the text of the preceding two sentences has been transmitted in a confused state.  It is most conspicuously observed in the non-sensical reading \emph{ārambhakriyārthatve tu sa tu} recorded in both \Tm\ and L\@.  The reading was already in the common archetype.  Even more curious is the reading \emph{vibhaktyarthaśravaṇam} in all the Devanagari manuscripts.  The scribes of \Td\ and M inserted \emph{°rtha°} when their respective exemplars did not have the syllable.  I do not believe that there was communication between \Td\ and M\@.  What probably happened was that two scribes independently and subconsciously inserted the syllable \emph{°rtha°} to form a compound that sounded more natural to them.}

In%\pratika{tasmād atha°}
\pratikap{p01_33} conclusion,\refm{31} it is reasonable that
the word \emph{atha\/} only has the meaning of announcing [a topic].

The%\pratika{`athe'ty a\-dhi\-kā\-rā°}
\pratikap{p01_34} meaning\refm{32} of the word \emph{iti} in
the ``the meaning of \atha{} is the declaration of the topic (\emph{athety
  adhikārārthaḥ})'' is to show that [the word \atha{}] is being quoted.  For
example, ``He said `cow' (\emph{`gaur' ity āha}).''  Even though the meaning of the word \atha{} is well-known, it is mentioned [by the \Bhasyakara] in order to make a student pay attention to what is going to be started.  Indeed the following formula
is well-known in \emph{śruti}:
\begin{quote}
        I am going to explain to you.  Pay close attention while I am 
        explaining.  (\rBAU{2.4.04/4.5.5}{2.4.4/4.5.5})
\end{quote}

Being%\pratika{kaḥ punar}
\pratikap{p01_35}\labelh{KHPNR1} asked,\refm{33} ``Then what is yoga,
whose instruction is announced [as the topic]?'' [the \Bhasyakara] says---\bht{Yoga is \samadhi.} %\footnote{Note that the author of the YVi does not
%  comment on the sentence:
%\begin{quote}
%        \emph{yogānuśāsanam adhikṛtaṃ veditavyam\dd}
%\end{quote}
%which appears in vulgate YBh.\label{note:veditavyam}}
 Since [the \Bhasyakara\ uses] the word `\samadhi\/ (assimilation),'\footnote{\label{ASSMLT}The translation ``assimilation'' for the word \samadhi\ is chosen here because usual translations of the word, such as ``concentration,'' etc., do not fit in the discussions below.  In the following \samadhi\ is said to be ``a property of the mind common to all stages.''  It is explicitly mentioned in our text that \samadhi\ is the state of stability in all the five stages of the mind.  See the paragraph starting with \emph{nanu ca bhūmīnām\ldots} below \reftoeditionpar{p01_42} (translated on p.~\pageref{NNCBH1}).} [he] rejects the meaning ``to unite'' of the root \emph{yuj-}, which [appears in the \Dhatupatha] as \emph{yujir yoge} (7.7), instead, he accepts the meaning ``assimilation (\emph{samādhi}),'' [which appears in the \Dhatupatha\ as] \emph{yuja samādhau} (4.68).  Yoga is assimilating (\emph{samādhāna}).\footnote{It is interesting to note that the two entries in the \Dhatupatha\ appear to reflect the two major trends in defining yoga.  One is, as seen here, to define it as \samadhi\ (perhaps originally meant as concentration, but not necessarily as describing a mental process).  The other is to define it as \samyoga.  One example of the latter definition (from the \Vaisesikasutra) is dealt with later in our commentary.  See pp.~\pageref{samyoga}\,ff. The same definition was presupposed by the \Naiyayika s (NS/NBh 4.2.38); even the \PYS\ 2.53--55 may have had the same system in its background.  The system of that yoga probably had the \Nyaya/\Vaisesika\ type of view on how perception occurs: by the contact between objects, sense faculties, \manas\ and \atman.  Their yoga was a gradual process to detach the component of perception from the outermost one, viz., first the objects, then the sense faculties.  When \manas\ resides with \atman\ (\samyoga), it is yoga.  See also \Carakasamhita\ \Sarirasthana\ 2.138--139.}
%These two \sutra s from the \emph{Dhātupāṭha} may
%  be of importance in studying the history of yoga, as according to
%  Bronkhorst~\cite{bronkhorst:history}, they entered the \emph{
%    Dātupāṭha} after \Patanjali.  This fight between whether \emph{
%    yoga} is \emph{samādhi} or \emph{saṃyoga} is reflected in the
%  discussion below. (Pages \pageref{samyoga}ff.)}

[Objection]\refm{34}%\pratika{nanu ca}
\pratikap{p01_36} If the \sutra\ is explained up to ``yoga is
\samadhi,'' then the \bhasya\ that starts with ``That (\samadhi) is a [property of mind] common to all stages'' is seen as detached.\footnote{This question arises only from the assumption that yoga/\samadhi\ is something special, that should be achieved with efforts.  The statement that it is present in all the stages of mind appears betrays such an assumption and thus the following discussions on the stages of mind appears to be detached. This is also the reason I have difficulty translating the word \samadhi\ in this context.  (See note \ref{ASSMLT} above.) Thus I think this question was posed by our author himself, as a hermeneutical problem.  He felt that the transition was abrupt, but was compelled to defend the text he was commenting upon.}

[Answer]\refm{35}%\pratika{naiṣa}
\pratikap{p01_37} There is no such fault.  For, when it is said that [yoga is] \samadhi\/ here, since [\samadhi] presupposes something to be assimilated[, i.e., the subject,] and since there are many things that can be assimilated, there is a dispute among those who know much,%
  \footnote{As the editors of the 1952 edition emended to \emph{bahuvidhā pratipatteḥ}, the expression \emph{bahuvidāṃ vipratipatteḥ} might seem awkward.  However, the construction, genitive plural plus the word \emph{vipratipatti} is found in \Sankara's works.  See \rBSBh{2.1.11}{2.1.11}: \emph{prasiddha\-māhātmyā\-bhimatānām api tīrtha\-karāṇāṃ kapila\-kaṇabhu\-kprabhṛtī\-nāṃ paraspara\-vipratipatti\-darśanāt}; 21: \emph{tatraivaṃ sati samyagjñāne puruṣāṇāṃ vipratipattir anupapannā}; \rBAUBh{4.5.15}{4.5.15}: \emph{vipratipattiś ca śāstrārthapratipattṝṇāṃ bahuvidām api}.  Cf.\ also \rBSBh{2.3.18}{2.3.18}: \emph{sa kiṃ kāṇabhujānām ivāgantukacaitanyaḥ, svato 'cetanaḥ, āhosvit sāṃkhyānām iva nityacaitanyasvarūpa eva, iti vādipratipatteḥ saṃśayaḥ}.  Here the word \emph{bahuvid} is probably used somewhat sarcastically.%

  The alternatives (\emph{ātman}, \emph{śarīra}, \emph{indriya\/}s) listed as the subject of \emph{samādhi\/} call for some attention.  These three, the most conspicuously the term \atman, are not part of the \PYS's terminology (an indication that our author is thinking of a wide context).  This set of three appears to reflect our author's view of what constitute an individual.  What is missing here is the word \manas\ or \citta. (Our author does not distinguish \manas\ and \citta.  He does not clearly distinguish \manas/\citta\ and \buddhi, either.)  But he omitted it because that is the correct answer to the question what is assimilated.  The set of \atman, \manas\ (or \citta), \sarira\ (or \deha), and \indriya s is what our author generally thought of as the constituents of an individual.  This is a rather simple view, somewhat reminiscent of the \Naiyayika/\Vaisesika\ view of an individual.%
  } %
such as ``Is \atman\ assimilated?''\ ``Is the body assimilated?''\ or ``Are the senses assimilated?'' Also, since assimilation presupposes something that qualifies itself [as in ``what exactly is it?'']%
  \footnote{See the next paragraph for what is considered to be the qualifier (\emph{viśeṣaṇa}) of \samadhi.%
  } %
[it is appropriate that the \bhasya\ continues the discussion on \samadhi, and hence there is no fault that the \bhasya\ up to \emph{yogaḥ samādhiḥ} and the following are detached].%
  \footnote{The translation is based on the understanding that the sentence that starts with \emph{iha samādhir ity ukte\ldots} consists of two main clauses even though it has four phrases in the ablative case expressing reasons.  This understanding derives from the positions of the particle \emph{ca}.  The first two phrases that end in the ablative case are combined by the particle \emph{ca}.  And semantically those are two stages.  This should lead to the first major reason \emph{bahuvidāṃ vipratipatteḥ}.  The second reason \emph{svaviśeṣaṇakākāṅkṣatvāc ca samādheḥ} is again connected with first major reason by the particle \emph{ca}.  This understanding fits the following commentary as well.  The two components of the question \emph{kasyāyaṃ, kiṃviśeṣaṇakaḥ} at the beginning of the next paragraph refer to the two reasons of the current sentence.  The first (\emph{kasyāyam}) corresponds to the question of what is assimilated, and the second \emph{kiṃviśaṣaṇakaḥ} corresponds to similarly worded \emph{svaviśeṣaṇkākāṅkṣatvāc ca samādheḥ}.%
  }

In\refm{36}%\pratika{kasyāyaṃ}
\pratikap{p01_38} order to answer the inevitable question, ``To what does
it (\samadhi) belong, what qualifies it?''\ [the \Bhasyakara] states the following---\bht{And that is a property of the mind common to all stages.}  It is a property of mind, not of \atman, etc.  And the mind assimilates itself because it does not require another agent.%

Being\refm{37}%\pratika{kāḥ punar}
\pratikap{p01_39} asked, ``Then what are those stages?''\ [the \Bhasyakara] answers---\bht{The stages are fixated (\emph{kṣipta}) [mind], stupefied (\emph{mūḍha}) % diff. fr. suggested
[mind]}, etc.  The past passive suffix \emph{-ta} in `fixated (\emph{kṣipta}),' etc., is used in the meaning that the object of an action is the agent itself (\emph{karmakartṛ}).  Compare, ``The granary collapsed by itself.''%
  \footnote{See MBh 3.1.87, \citet[nn.~22--24]{wezler:yvi1}\index{Wezler, Albrecht}.%
  }  %
\bht{Fixated [mind]} means a benumbed [mind] being fixed with an undesirable object.  \bht{Stupefied [mind]} lacks the distinction [between different objects].  \bht{Scattered (\emph{vikṣipta}) [mind]} means being fixated at several objects [where the passive past participle is used] in the same sense [as in the case of fixated mind] in which the object of the action is the agent itself.  Also, it is not capable of distinguishing [different objects] for the very reason that it is scattered.  \bht{Focused [mind] (\emph{ekāgra})} means the [mind] holding a consistent notion.  \bht{Suppressed [mind] (\emph{niruddha})} is without any notion---[All the adjectives modify the word] the mind.%

[Objection]\refm{38}%\pratika{nanu ca}
\pratikap{p01_40} If [the \Bhasyakara] intends to talk about stages [of the mind], properties (\emph{dharma}), why is an property-bearer \emph{dharmin}) [that is the mind] discussed by ``fixated [mind],'' etc.?%
  \footnote{Note that all the terms to refer to the different stages of mind in the \bhasya\ are adjectives in the nominative case, neuter singular.  They naturally modify the word \emph{cittam}.  And this is the understanding our author holds and is the background of this objection.%
  }

[Answer]%\pratika{naiṣa}
\pratikap{p01_41}\refm{39} There is no such fault.  The property-bearer indicates exactly the property.  Because the properties belong to the property-bearer.  For
example: If one is asked ``What is the characteristic of being a
cow?''\ the property is indicated by the property-bearer, by saying, ``It has
horns and a hump, and is hairy toward its tail.''\footnote{According to
  \citet{wezler:yvi3}\index{Wezler, Albrecht}, this is a reference to \rVS{2.1.18}.  Wezler
  points out that the \sutra\ is referred to in support of the
  view held by the author of the YVi, and that the \emph{sūtra}, therefore,
  was favorable to him.  Of course the archetypal description of the cow appears already in the \VMBh.} Therefore [the \Bhasyakara] means to say that fixation, etc., are the stages, the properties of the mind.

% nanu ca bhuumiinaa.m 
[Objection]%\pratika{nanu ca}
\pratikap{p01_42}\labelh{NNCBH1}\refm{40} If the stages are properties of the mind and \samadhi, too, is [a property of the mind], then how do we distinguish \samadhi\/ that is said to be ``common to all the stages'' as a common denominator from the stages?  [If there is nothing special about \samadhi, it would not be necessary to introduce the term.]%
     \footnote{The whole exchange appears a little clumsy.  Answering to the question ``How (\emph{ka\-tha\-m})\ldots?'' with ``no'' does not appear to be usual.  I have interpreted and translated it as short for ``naiṣa doṣaḥ'' below.  Still, some text might have been lost.%
     }
%
% na| saamaanyabhuuta
[Answer]\refm{41} That is not correct.  Because \samadhi\/ is
universal, while the stages are particulars.  That is, in such phrases as ``[the mind] stays fixated,'' ``[the mind] stays stupefied,'' ``[the mind] stays scattered,'' and ``[the mind] stays focused,'' and ``[the mind] stays suppressed,''%
  \footnote{I have conjecturally added \emph{niruddham} after \emph{vikṣiptam} and \emph{ekāgram}, while it is missing in all the manuscript and the \EP.  Also, following the \EP, I have emended \emph{vikṣiptam} at the beginning of the series to \emph{kṣiptam} and \emph{kṣiptam} at the third position to \emph{vikṣiptam} so that the order of the stages follows that of the \bhasya.  I do not believe that the swapped positions of \emph{kṣipta} and \emph{vikṣipta} or the lack of \emph{niruddha} were intended by the author.  I judge rather that those are results of the damage to the text during the transmission.  Cf.\ the following for the state of the text preserved in manuscripts.  Frequent appearances of the same words (\emph{kṣiptam}, etc.) may have contributed to the corruption in this particular sentence.%
  } %
stasis (\emph{sthiti}) in the stages ``fixated,'' etc., accompanies universally.  And\labelh{station_immersion} \samadhi\/ is stasis.  Therefore [\samadhi] is commonly%
  \footnote{The reading \emph{sā\-dhā\-ra\-ṇye\-na} is a result of emendation.  Most manuscripts read \emph{prā\-dhā\-nye\-na}, except \Tm\ \emph{sā\-dhā\-nye\-na} before correction.  It is likely that the common archetype of \Tm\ and L indeed read \emph{prā\-dhā\-nye\-na}.  What probably happened was that the scribe of \Tm\ first wrote \emph{sā\-dhā\-nye\-na}, expecting the word \emph{sā\-dhā\-ra\-ṇye\-na} before realizing that he had written something not in his exemplar.  Nonetheless, I adopt the reading what he subconsciously expected; otherwise, the sentence is not coherent.%
  } %
present in all the stages in the form of a universal.%
  \footnote{The word \emph{iti} at the end of this paragraph marks the end of a topic, and hence no specific translation is given.%
  }

% sarvabhuumip.rthivii
% Vetter's trl.:
[According to \rAAP{5.1.41--42}{5.1.41--42}]%\pratika{`sarvabhūmi°}
\pratikap{p01_43} \emph{sarvabhūmipṛthivī}[\emph{bhyām aṇañau} and \emph{tasyeśvaraḥ}] the suffix \emph{a} (\emph{aṇ}) [added to \emph{sarvabhūmi}] gives the meaning ``governing [all stages].''  Since the rule \emph{anuśatikādī\emph{[}nāñ ca\emph{]}} (\rAAP{7.3.20}{7.3.20}) is applicable [because \emph{sarvabhūmi} is included in the \emph{gaṇa} \emph{anuśatikādi}], [like in \emph{ānuśātika}, the first vowel] in both constituents of the compound [\emph{sarvabhūmi}] is substituted by the longest modification[, thus resulting in the form \emph{sārvabhauma} as found in the \bhasya].

%[The\refm{42} word \emph{sārvabhauma} has] the \emph{aṇ} suffix
%(\emph{a}), because it is governed by [P\=a\d{n}ini's rule] ``\emph{%
%  sarvabhūmipṛthivī}'' (AA 5.1.41).  [The \emph{aṇ} is used] in
%the meaning of `governing'[, the meaning mentioned in the subsequent
%rule ``\emph{tasyeśvaraḥ}'' (AA 5.1.42)].  Thereupon the resulting
%compound is governed by the rule ``\emph{anuśatikādīnāñ ca}''
%(AA 7.3.20), resulting in the \emph{vṛddhi} [substituting the first
%vowel] of both constituents of the compound.

% anye punar
On%\pratika{anye}
\pratikap{p01_44}\refm{43}\labelh{ANYEPUNARtrl} the other hand, some people explain the stages as external and internal objects of \samyama\ meditation.%
  \footnote{The relation of \samyama\ with particular objects and its result is the main topic of the third {pāda}.  Here the view that the \emph{bhūmi\/}s in the context of \rYBh{1.01}{1.1} is the same \emph{bhūmi\/}s in the context of \emph{saṃyama\/}s in \rYS{3.06}{3.6} (\emph{tasya bhūmiṣu viniyogaḥ}) is criticized.  Someone held the view (according to the YVi) that they are the same, presumably because the same word \emph{bhūmi\/} is used in both places.  It should be noted that the author of the YVi is referring to an interpretation of the YBh here.  It is probable that he is using a sub-commentary on the YBh.  Note, however, that the interpretation is different from that of \Vacaspati.%  See section~\ref{vpmandyvi}.

  Philologically, our manuscripts indicate that their common ancestor consistently read \emph{samayama} instead of \samyama.  It is interesting that such an obviously faulty reading was consistently adopted and more or less faithfully reproduced.%
  } %
[But] for them there is no chance to place [a stage and the mind] in the same case (\emph{sāmānādhikaraṇya} relationship) [in the way the \Bhasyakara\ does:] ``When the mind is scattered\ldots\@,''%
  \footnote{This is a reference to the next sentence in the \bhasya.  See page~\pageref{ATRVKSP}\,ff.%
  } %
and they contradict their own doctrine.  How?  [According to their interpretation,] that on which [the mind] is fixated is ``fixated''; that with regard to which  [the mind is stupefied] is ``stupefied''; and that with regard to which [the mind is scattered] is ``scattered.''%
  \footnote{The last phrase \emph{vikṣpiyate 'sminn iti ca vikṣiptam} is missing in all the manuscripts and the \EP.  The context requires this last item to be mentioned as well.  One, albeit not very strong, indication that there is a textual corruption is the lack of the particle \emph{ca} in the preceding phrase regarding \emph{mūḍha}.%
  }  %
If this [interpretation] were true, then they could not be the object of \samyama\ meditation.  In that (in the state of \samyama) \emph{kṣepa}, etc., cannot occur.  For, \samyama\ meditation is possible at the state of focused (\emph{ekāgra}).

Furthermore,%\pratika{kiñ cānyat}
\pratikap{p01_45}\refm{44} since the root \emph{kṣip-} does not have the
meaning of ``steadiness,'' etc., it is not possible to have the \emph{%
  niṣṭhā} suffix in the meaning of locus (\emph{adhikaraṇa}) [as prescribed by {%
  \Panini}].\footnote{See \rAAP{3.4.76}{3.4.76}: \emph{kto ’dhikaraṇe ca dhrauvyagatipratyavasānārthebhyaḥ} (The past passive participle affix \emph{-ta\/} is used in the sense of the locus, too, when it follows verbal roots that have the meanings of steadiness, motion, or termination).  The argument here is still based on the analysis of the opponent's view: \emph{kṣipyate ’sminn iti kṣiptam}, etc.  Note the use of the locative case (\emph{yasmin}).  It represents the locus (\emph{adhikaraṇa}).  Our authors angle of attack here is that such a formulation, ``that on which something is thrown is `thrown'\thinspace,'' to analyze a past passive participle (\emph{kṣipta}, etc.)\ is grammatically impossible.  Pāṇini allows only verbs that have certain meanings to form a passive past participle in the meaning of the locus.  The root \emph{kṣip-} does not have any of the three meanings, \emph{dhrauvya}, \emph{gati}, or \emph{pratyavasāna}.  According to the \Dhatupatha, the root \emph{kṣip-} has the meaning \emph{preraṇa} (driving).}

Furthermore,%\pratika{kiñ ca}
\pratikap{p01_46}\refm{45} when the mind has been suppressed, \emph{saṃyama\/} meditation, too, will not exist because no activity exists.  Accordingly, the object of \samyama\ meditation does not exist.  For, if suppression was a specific object, the mind would not be suppressed.  Because at that time (when the mind is suppressed) [the mind] ceases to take an object.  Also, [there is a problem that] completely enumerating [the stages] is not possible.  Because the objects of \emph{saṃyama\/} meditation are not just five of \emph{kṣipta}, etc.  Because they are innumerable.

Objection:%\pratika{nanu ca}
\pratikap{p01_47}\labelh{samyoga} Some\refm{46} consider yoga differently.  Namely, they say---\footnote{\citet{wezler:yviandvs}\index{Wezler, Albrecht} deals with this
quotation from/reference to the \VS\ in detail.}
\begin{quotation}
        \bht{Pleasure%\pratika{indriya}
\pratikap{p01_48} and pain arise from the contact between senses, mind, and objects.  [Yoga] is a union (\emph{saṃyoga}),}%
		\footnote{All the manuscripts read \emph{saṃyogaḥ} at the end of the second \sutra.  \citet{wezler:yviandvs}\index{Wezler, Albrecht} proposes to emend this to \emph{sa yogaḥ}.  But this emendation does not seem necessary.%

		The commentary that is cited with these \sutra s below does not support \emph{sa yogaḥ} in the \sutra.  First, the word \emph{sa} in the commentary \emph{sa prāṇamanovinigrahāpekṣaḥ} towards the end of the quoted commentary corresponds to the preceding relative pronoun \emph{yaḥ} in \emph{yo ’sau}.  The syntax of the commentary should be viewed as ``\emph{yo ’sau\ldots saṃyogaḥ} (a word from the \sutra) \emph{sa prāṇamanovinigrahāpekṣaḥ} (a word from the \sutra) \ldots.''  Accordingly, the word \emph{sa} does not seem to be from the \sutra.  The parallel syntax in ``\emph{yo ’sau} \ldots \emph{sannikarṣaḥ} \ldots \emph{tasya}'' near the beginning of the commentary supports this view.  The commentary employs the similar syntax when explaining the relationship between words in the \sutra s. %
		Secondly, the commentary is explicit in that yoga is \emph{saṃyoga\/} and treats \emph{sukhaduḥkhābhāvaḥ} more like a condition (\emph{tasyām avasthāyām}).  [Cf.\ the Mithilā version of \rVS{5.2.14}: \emph{saśarīrasya sukhaduḥkhābhāvaḥ\dd} \citep[113]{vs} where the \sutra\ is terminated after \emph{sukhaduḥkhābhāvaḥ} and the word is not syntactically connected with the following \sutra: \emph{saṃyogaḥ\dd}; and \rnnNS{4.2.45}{4.2.45}/\rnnNBh{4.2.45} 4.2.45: \emph{\textbf{tadabhāvaś cāpavarge\ddd} tasya bu\-ddhi\-ni\-mittā\-śraya\-sya śa\-rīre\-ndri\-ya\-sya dha\-rmā\-dharmā\-bhāvād a\-bhā\-vo ’pa\-va\-rge\dd\ ta\-tra yad u\-kta\-m ``apa\-varge ’py evaṃ pra\-saṅga'' iti (4.2.43), tad a\-yu\-ktam\dd\ tasmāt sarva\-duḥkha\-vi\-mokṣo 'pa\-vargaḥ\dd\ yasmāt sarva\-duḥkha\-bījaṃ sa\-rva\-duḥkhā\-ya\-ta\-naṃ cā\-pa\-varge vi\-cchi\-dyate tasmāt sa\-rve\-ṇa duḥkhe\-na vi\-mu\-ktir apa\-vargaḥ\dd\ na nirbī\-jaṃ ni\-rā\-yatanaṃ ca duḥ\-kham utpa\-dya\-ta iti\ddd}.]  This makes it difficult for \emph{sukhaduḥkhābhāvaḥ} and \emph{yogaḥ} to be in the \emph{sāmānādhikaraṇya\/} relationship, at least according to the commentary. %
		Thirdly, the commentary refers to \emph{saṃyoga\/} twice. This insistance on \emph{yoga = saṃyoga\/} makes it easier to think that the word \emph{saṃyogaḥ} was already in the \sutra\ as part of the definition of yoga than that it was introduced from somewhere else. %
		Thus even if the word \emph{saṃyogaḥ} might not be original to the \VS, there seems to be little support that the author of the YVi \emph{did not} quote from the \VS\ that already contained the word.%

		Apart from the commentary quoted in the YVi below, evidence from the Vaiśeṣika side appears to support the presence of \emph{saṃyogaḥ\/} in the \sutra, too.  All the manuscripts of the VS used by Jambūvijaya have \emph{saṃyogaḥ} \citep[42]{vs}.%

		In addition, there appears to have been a great deal of controversy over the definition of yoga---whether it is \emph{saṃyoga\/} or \emph{samādhi}---before the authority of the YS/YBh prevailed.  This may be reflected even in two entries of the root \emph{yuj-\/} in the Dhātūpāṭha, namely \emph{yujir yoge} and \emph{yuja samādhau}.  The two entries of the Dhātupāṭha are said to be referred to in the \YBh. (See the paragraph \emph{kaḥ punar yogo\ldots} \reftoeditionpar{p01_35} translated on p.~\pageref{KHPNR1}.  Cf.\ \citet{bronkhorst:history}.)  Another reflection of this controversy may be seen between \rNS{4.2.38}{4.2.38} (\emph{samādhiviśeṣābhyāsāt}) and its interpretation in the NBh (\emph{sa tu pratyāhṛtasyendriyebhyo manaso dārakeṇa prayatnena dhāryamāṇasyātmanā saṃyogas tattvabubhutsāviśiṣṭaḥ} \ldots).  The \NBh\ uses \emph{saṃyoga\/} instead of \emph{samādhi\/} of the \sutra.  [The system of yoga adopted in the Nyāya and Vaiśeṣika systems is closely tied to their epistemology.  Cf.\ the criticism toward the Nyāya/Vaiśeṣika theory of the generation of knowledge through connection (\emph{saṃyoga}) between \emph{manas\/} and \emph{ātman} in the YVi on \PYS\ 1.2.% (\pageref{nvsamyoga}--\pageref{nvsamyogaend}).%  Close relation between yoga and the knowledge of the truth \emph{tattvajñāna\/} may also be reminded.  In the yoga system adopted in Nyāya/Vaiśeṣika systems, the goal of yoga is \emph{tattvajñāna\/} while in the system of Patañjali the goal is \emph{cittavṛttinirodha}.%
		]  Thus given the conflict between \emph{saṃyoga\/} and \emph{samādhi\/} surrounding the definition of yoga, it seems natural that a sūtra defining yoga included the other synonym for yoga.%

		And finally the criticism in the paragraph \emph{athāpi syāt---yogaḥ samādhir iti\ldots} \reftoeditionpar{p01_63} (translated on p.~\pageref{ATHPSY1}) by the author of the YVi presupposes a defining term (\emph{samādhi\/} or \emph{saṃyoga}) to be present in the \sutra.  Note that when he finally deals with the view \emph{yogaḥ samādhiḥ}, he has criticized most of the words in the opponent's sūtras:  \emph{indriyamano'rthasannikarṣāt} in two paragraphs \emph{kiñ cānyat\dd\ mano°\ldots} \reftoeditionpar{p01_55} (trl.\ p.~\pageref{KCMN1}); \emph{sukhaduḥkhe} in three paragraphs starting with \emph{kiñ cānyat---sa\-rva\-prāṇa°\ldots} \reftoeditionpar{p01_57} (trl.\ p.~\pageref{SRVPRNBH1}); \emph{tadanārambhaḥ} in the paragraph \emph{indriyasaṃyoge\ldots} \reftoeditionpar{p01_62} (trl.\ p.~\pageref{INDRYSMYG1}); \emph{ātmasthe manasi} in four paragraphs starting with these exact words \reftoeditionpar{p01_51} (trl.\ p.~\pageref{ATMSTH1}); \emph{saśarīrasya} in the paragraph \emph{kiñ cānyat---upātte 'pi\ldots} \reftoeditionpar{p01_60} (trl.\ p.~\pageref{UPTTP1}); \emph{sukhaduḥkhābhāvaḥ} in three paragraphs starting with \emph{kiñ cānyat---sarvaprāṇa°} \reftoeditionpar{p01_57} (p.~\pageref{SRVPRNBH1}); \emph{prāṇamanovinigrahāpekṣaḥ} in the paragraph that starts with \emph{kiñ cānyat---prāṇamano°\ldots} \reftoeditionpar{p01_61} (trl.\ p.~\pageref{KCPRMN1}).  There should be something to be criticized in the paragraph \emph{athāpi syāt---yogaḥ samādhir iti\ldots}.  Note again that in \rNBh{4.2.38} \emph{samādhi\/} in the \sutra\ is replaced by \emph{saṃyoga}.%

		From the above, it appears that there was the word \emph{saṃyogaḥ} in the \sutra\ in question.  Adding \emph{yogaḥ}, as suggested by the editors of the 1952 edition (p.~6, l.~10) might make the \sutra s complete as a definition of yoga.  The commentary of Candrānanda also seems to imply that there was the word \emph{yogaḥ} as the last element of the \sutra\ \citep[42,24]{vs}.  Despite my inclinations, both from those of the YVi and the \VS, I did not add the word in the \sutra.  The fact that the author of the YVi supplies the word \emph{yogaḥ} to both \emph{sukhaduḥkhābhāvaḥ} (here is a deviation from the interpretation given in the commentary accompanying the \sutra s) and to \emph{samādhiḥ} (the replacement I think for \emph{saṃyogaḥ}) suggests that the word \emph{yogaḥ} was provided by the context of surrounding other \sutra s.%  After all, the \sutra s in question are viewed as not original to the \Vaisesikasutra\ \citep{wezler:yviandvs}.  Perhaps we should not forget that there is even a possibility of interpolation (in the VS) and quotation (in the YVi) from the same source.
		} %
	\bht{which follows restraining of breaths and mind, and is [characterized as] its (of the contact) not happening, absence of pleasure and pain, and it belongs to someone who still has a body, under the condition that the mind stays with \atman.}%
        
    The%\pratika{tad iti}
\pratikap{p01_49} word ``\bht{Its}'' presupposes something that is being        discussed---the cause \bht{of pleasure and pain} = the \bht{contact between} \atman, \bht{senses, mind, and objects}.  Its \bht{not happening} = its not originating.  How does it occur?%
        \bht{Under the condition that the mind stays with \atman}---not when it stays with senses.  [Yoga] \bht{belongs to someone who still has body}---belongs to someone whose body has not perished. At that time, as the contact does not exist, \bht{pleasure and pain} do \bht{not exist}.   For it is said ``if the cause does not exist, its result does not exist (\rVS{1.2.1}).''  Under this circumstance there is a \bht{union} between omnipresent \bht{\atman} and the \bht{mind}.  That is yoga, which is a specific \bht{union} that \bht{follows restraining of breaths and the mind}.%
\end{quotation}

We%\pratika{atrocyate}
\pratikap{p01_50} answer---

[Answer]%\pratika{`ātmasthe}
\pratikap{p01_51}\labelh{ATMSTH1} To say\refm{47} that \bht{under the condition that the mind stays with \atman} is not rational.  Because the mind always stays with \atman.
  
If%\pratika{indriyādi°}
\pratikap{p01_52} you say, ``[It is said] with regards to [the condition] that contact between senses, etc., is not happening,'' then stating ``under the condition that the mind
stays with \atman'' would be meaningless since it is already established
just by saying ``its not happening.''

Furthermore,%\pratika{kiñ cānyat}
\pratikap{p01_53} even a liberated (\emph{mukta}) would have yoga since his mind stays with \atman, and he does not have contact with senses.\footnote{Here our author assumes that his opponent would not like the idea that a liberated person has yoga.  If the opponent holds the view that yoga is the direct means of liberation, this might not be a problem. This does not appear to be the case.  Liberation (\moksa) is defined in the Vaiśeṣikasūtra a few sūtras after the sūtras under consideration here: \emph{ta\-da\-bhā\-ve saṃyogābhāvo ’prādurbhāvaḥ sa mokṣaḥ} (5.2.20 in the numbering with Candrānanda's commentary).  Hence it is indeed probably undesirable that a liberated person has yoga.  It is, however, interesting to note that our author assumes that the ``liberated'' still has \manas.  The \emph{saṃyoga} in \emph{saṃyogābhāvaḥ} in the above sūtra is, according to Candrānanda, the union between \atman\ and \manas, and this interpretation is the most plausible.  (The word \emph{tad} in \emph{tadābhāve} refers to \emph{adṛṣṭa}s mentioned in the previous sūtra.)  Our author thus has a different understanding of what a ``liberated'' is.  The expression \emph{muktātman} or simply \emph{mukta} is used several times in this work, especially in relation to \Isvara.  For example, see p.~\pageref{LBRTDATM}.  There the assumption is that the liberated ātman or a liberated person does not have a body.  The use of the term is probably consistent with that of Kumārila.  See note \ref{KMRSAP} on p.~\pageref{KMRSAP}.}  For, since he, too, [being \atman,] is omnipresent, and the mind is eternal,\footnote{The \Vaisesika s hold that \manas\ is eternal.  See the references in the edition part.} [the mind] would indeed always stay with \atman.
% 
% Usage of `mukta' seems to be interesting because he is speaking of 
% a special being.
% See PDhS and p.84 (370):
% \begin{quotation}
%        yathāsmadādīnāṃ gavādiṣvaśvādibhyas
%       tulyākṛtiguṇakriyāvayavasaṃyoganimittā
%       pratyayavyāvṛttir dṛṣṭā gauḥ śuklaḥ śīghragatiḥ
%       pīnakakudmān mahāgaṇṭa iti. tathāsmadviśiṣṭānāṃ
%       yogināṃ nityeṣu tulyākṛtiguṇakariyeṣu paramāṇuṣu
%       muktātmamanaḥsu cānyanimittāsambhavād yebhyo
%       ^^^^^^^^^^^^^^^^^
%       nimittebhyaḥ pratyādhāraṃ vilakṣaṇo 'yaṃ vilakṣaṇo
%       'yam iti pratyayavyāvṛttiḥ, deśakālaviprakarṣe ca
%       paramāṇau sa evāyam iti pratyabhijñānaṃ ca bhavati, te
%       'tyā viśeṣāḥ.
% \end{quotation} , also cf. p.80 (359), p.84 (371)
% 
% And by the way rememver dāruyantra in Vivaraṇa?  PDhS has an
% entry---i.e. śarīraparigṛhīte vāyau vikṛtakarmadarśanād
% bhastrādhimāpayiteva.  nimeṣonmeṣakarmaṇā niyatena
% dāruyantraprayokteva.

Furthermore,%\pratika{kiñ cātmanaḥ}
\pratikap{p01_54}\refm{49} since there are no regions in \atman,\footnote{That \atman\ has no region is not part of the \Vaisesika\ system. Old \Vaisesika s take for granted that \atman\ has regions. This is observed in expressions such as \emph{antarhṛ\-da\-ye ni\-ri\-ndri\-ye ā\-tma\-pra\-de\-śe ni\-śca\-lam ma\-nas ti\-ṣṭha\-ti} \citep[183]{pdhs}, \emph{ha\-sta\-vaty ā\-tma\-pra\-de\-śe pra\-yatnaḥ saṃjā\-ya\-te} \citep[297]{pdhs}. Note that the expression \atmapradesa\ is used in the sense of “a region in \atman.” The \Naiyayika s, on the other hand, saw a problem in expressions such as \emph{ātmapradeśa}. \Nyayasutra\ \rnNS{2.2.17}{2.2.17} reads \emph{kāra\-ṇa\-dra\-vya\-sya pra\-deśa\-śabde\-nā\-bhi\-dhā\-nāt}. This is an answer to the criticism expressed in \rNS{2.2.14}{2.2.14} \emph{na, ghaṭā\-bhā\-va\-sā\-mā\-nya\-ni\-tya\-tvān ni\-tye\-ṣv apy a\-ni\-tya\-vad upa\-cā\-rāc ca}. The part \rNS{2.2.17}{2.2.17} is responding to is \emph{nityeṣv apy anityavad upacārāt}. This, in turn, is an objection to one of the reasons that the proponent gives to establish that \sabda\ (sound) is impermanent (\anitya) and to be created (\emph{utpadyate}) rather than revealed (\emph{abhivyajyate}). It is the part \emph{krṭakavad upacārāt} in \rNS{2.2.13}{2.2.13}. Following the explanation of the \NBh, this part of the discourse goes as follows: the sound is impermanent since like any other things that are made (\krtaka), i.e., impermanent, the sound gets the figurative expression such as intense or mild (2.2.13). The opponent argues that the reasoning has a deviation (\vyabhicara) since eternal things also get figurative expressions in the same ways as the things that are made do. Examples are a region of space (\akasa) or a region of \atman\ (2.2.14). The assumption is that only the things that are made (\krtaka) have parts, viz., regions, and they can be referred to. But we refer to regions of eternal, viz., non-made, things---presumably without regions---figuratively. Then the proponent replies that in the example the opponent gives, a region of eternal things (\akasa\ or \atman), the referent is not the eternal things but the fact that the union between the eternal thing and [spatially] limited substance does not fill the whole eternal thing (\emph{saṃyogasyāvyāpyavṛttitvam}). That is, the opponent's example does not annul the argument that the sound gets figurative expression just like any other things that are made. For, the expression that purportedly show eternal things being referred to in the same manner as impermanent things in fact does not refer to the eternal things. Rather, it is only a derivative (\emph{bhākta}) understanding that \akasa\ or \atman\ has a region.

Admittedly, it is hard to precisely understand what is meant by the root \emph{upa-car-}, which should be understood as ``to figuratively express,'' in this discourse. However, we can still glean that the \Naiyayika s were concerned about the expressions such as \emph{ākāśapradeśa} or \emph{ātmapradeśa} in their system adopted from the \Vaisesika s where \akasa\ or \atman\ should not have a region. It is also possible that their solution to the problem was that those are figurative or secondary expressions. This is seen in the next sentence in our text.

It is interesting that our author states that \atman\ has no regions matter-of-factly when there is no statement to the effect in old \Vaisesika\ works. \Sankara\ assumes the same. For example, in \rBSBh{2.2.17}{2.2.17} he simply states that \anu s, \atman, and \manas\ have no regions (\emph{aṇvātmamanasām apradeśatvāt}). In \rBSBh{2.3.53}{2.3.53} the topic is a region of omnipresent \atman\ delimited by a body. There he argues that, even assuming multiple \atman s, since \atman s are equally omnipresent (following the \Vaisesika\ doctrine) all the \atman s are inside all the bodies and thus it is impossible to postulate a region of ātman, even when it is delimited by a body. There is a leap of logic here. As he further discusses, it may be impossible to discern which \atman\ among all the \atman s is being affected by a body, yet it does not follow that the part of \atman\ inside a body and the rest cannot be distinguished. He seems to believe that omnipresent means having no region. He further qualifies the word \emph{ātman} with the adjective \nispradesa, and for such \atman, a region delimited by a body is ``made up (\emph{kālpanika})'' and incapable of determining a real effect. Along with the assumption that there is no regions in \atman, the mention of ``made up'' region of having no effect is similar to the discussion here. For the text of the BSBh, see p.~\pageref{BSBHTHREETHREE}} to say  ``under the condition that the mind stays with \atman'' is not rational.  Also, even if the union with a region of \atman\ is said figuratively, it will not be the cause of true yoga.  For something figurative is [ultimately] a false.

There%\pratika{kiñ cā\-nya\-t\dd\ ma\-no°}
\pratikap{p01_55}\labelh{KCMN1}\refm{50} is another [difficulty in your view].  The use of the word ``senses (\emph{indriya})'' would be pointless since [that contact with objects is] already established just by saying ``from the contact between the mind and objects.''

If%\pratika{athāpi}
\pratikap{p01_56} you\refm{51} may say---[Objection] But since the mind and objects
cannot have contact without senses, it would be irrational to negate
something that cannot be achieved, saying ``its not happening.''

It\refm{52} is not so.  Without senses, the contact between the mind and objects
does not exist.  Thus naturally, the person suppressing the contact should already know that the acquisition of objects is through senses.

There%\pratika{kiñ cānyat---sarva\-prāṇa°}
\pratikap{p01_57}\labelh{SRVPRNBH1}\refm{53} is another [difficulty in your view].  Because it is well
known that anything that breathes has pleasure and pain from the
contact between the senses, the mind and objects, and because the connection between \atman\/ and the mind is eternal, [what is meant] would be established by
merely saying ``yoga is the absense of pleasure and pain.''  Therefore the
rest of the \sutra\ would be meaningless.

There%\pratika{athāpi syād}
\pratikap{p01_58}\refm{54} might be the following objection: [Objection] If the \sutra\ says as much as ``[yoga is] the absence of pleasure and pain'' without
mentioning the contact between senses, the mind and objects, then,
since a liberated person does not have pleasure and pain, there would be an undesirable consequence of he becoming a yogin.  Therefore in order to avoid this [undesirable consequence], the rest of the \sutra\ is included.

[Answer]%\pratika{tac ca na}
\pratikap{p01_59}\refm{55} That again is not correct.  Because a liberated 
does not acquire pleasure and pain.  [And] because it is logical that
negation is done on something that has previously been acquired.  Someone
acquires pleasure and pain, and yoga is only the absence of pleasure
and pain with regard to him.  Therefore it is reasonable to
reject the rest of the \sutra.

There%\pratika{kiñ cānyat---upātte 'pi}
\pratikap{p01_60}\labelh{UPTTP1}\refm{56} is another [difficulty in the \sutra].  Even though the
word ``[that] belongs to someone who still has a body (\emph{%
  saśarīrasya})'' is adopted, it is impossible to exclude a liberated (\emph{%
  mukta}).  Because it lacks the purpose.  If yoga is the absence of
pleasure and pain, then the body does not mean anything, even when
it still exists, because it is does not produce any effect.  When we do not take the effect into consideration, there is absolutely no distinction between a liberated and not liberated.
                                % who is mukta?

There%\pratika{kiñ cā\-nya\-t---prā\-ṇa\-ma\-no°}
\pratikap{p01_61}\labelh{KCPRMN1}\refm{57} is another [problem in the \sutra].  The part ``[which]
follows restraining of breaths and mind'' is not appropriate, either.
Because the effort that is holding the restraining is inherent to
\atman.  Even though the mind does have activities, without being
connected with senses, it is impossible to restrain the mind and
breaths.
                                % manasaḥ kriyāvattva?
                                % prayatnasya ātmasamavetatva?

Also,%\pratika{indriyasaṃyoge}
\pratikap{p01_62}\labelh{INDRYSMYG1} if\refm{58} the mind is connected to senses, it is not reasonable to
say ``its (contact's) not beginning.''  Also, since the mind is immersed in
restraining breaths, it is impossible for it to not begin [contacting
with senses].  Also, if there is no contact between the senses and
the wind, there is no restraining of breaths.  Also, if breaths are 
restrained, there is no yoga.  Because the cause of restraining breaths
is the relation with the action in the mind through the holding effort
residing in \atman.

Then%\pratika{athāpi syāt---yogaḥ samādhir iti}
\pratikap{p01_63}\labelh{ATHPSY1}\refm{59} you might say---[Objection] Yoga is \samadhi.  [Answer] %%\pratika{tacca na}
% \pratikap{p01_64} 
That is not right.  Because [saying so] does not have any effect.  Again, \atman\ and manas are eternally assimilated.  Also, it is said that \samadhi\/ is stasis.\footnote{See page \pageref{station_immersion}.}

In%\pratika{tasmād}
\pratikap{p01_65} conclusion\refm{60}, it is appropriate to say ``And that (\samadhi) is
a property of the mind common to all stages.''

[Objection]%\pratika{yady evaṃ}
\pratikap{p01_66}\refm{61} If, as you say, yoga is \samadhi, and it is a
property of mind common to all stages, then it is [already]
established for all the living creatures without effort, because [it]
includes fixated [mind], etc.  Accordingly, the definition (\emph{lakṣaṇa})
is meaningless because it is established without effort, just
like inhaling and exhaling.  Thus the instruction of yoga (\emph{yogānuśāsana}) also becomes meaningless.

Answering%\pratika{atra vikṣipte}
\pratikap{p01_67}\refm{62}\labelh{ATRVKSP} this, [the \Bhasyakara]
states---\bht{among those, if the mind is scattered, the \samadhi\/ that
  is dominated by being scattered does not belong to the side of yoga.}
For, [\samadhi] that includes ``fixated [mind],'' and so on, is not intended to be on the side of yoga.  Namely, \samadhi\/ that includes [the condition of] being fixated, etc., is not capable of illuminating an object as is.  Because the action of being fixated is predominant in it.  By merely negating [the stage of] scattered (\emph{vikṣipta}), according to the maxim of defeating the representing
wrestler, the stages fixated (\emph{kṣipta}) and stupefied ({mūḍha}) are also negated.  \emph{Samādhi\/} of being scattered (\emph{vikṣepasamādhi}) represents [other types of \samadhi\/ that do not belong to the side of yoga] because it [still] has an adequate qualification [to be included on the side of
yoga].  Namely, the mind which is in the state of scattered is able
to acquire desirable object, for it has not fallen into [a
particular] side.  But stupefied [mind]---by means of attaching [it] to an object, for example---or fixated [mind]---by means of detaching [it] from undesirable object, for example---cannot be led to anywhere else.  [The \Bhasyakara] says ``the side of yoga''---Even though it is still
\samadhi, since it does not produce an effect [expected from yoga], he says ``the side.''  For example---even though a walking man stops at every step, since it
does not cause the effect of stopping, it is not called
standing.\footnote{Cf. {Nāgārjuna}\/'s famous discussion on going/walking (\gati)
  in {Madhyamikakārikā}s \citep[kk.~2.1--25]{mmk}.  Note that
  while {Nāgārjuna} tries to prove there is no such thing as
  ``going,'' the point the author of the YVi making is the
  opposite.}
  
With%\pratika{`vikṣepo°}
\pratikap{p01_68}\refm{63} the intention [of specifying] a particular kind of \samadhi\/ that is in the state of being scattered, [the \Bhasyakara\ refers to] the
quality of being ``dominated by being scattered.''  Because, on the other hand, [that yoga is \samadhi] primarily applies to all the stages equally, for [the \Bhasyakara] has stated ``common to all stages.''%. Since it is stated For, on the other hand, the property of being common to all the stages main
%[characteristic of immersion] is that it is ``common to all stages''.
%[It is main characteristic] because it exists at all stages.

If%\pratika{yady upa°}
\pratikap{p01_69}\refm{64} one might question---If being dominated is the cause of not belonging to the side of yoga, now, in the case of focused [mind] there would be the dominance of being focused.  Then [the \samadhi\/ in the stage of focused] does not belong to the side of yoga.  In order to answer this, [the \Bhasyakara]
states---\bht{But [\samadhi] when [the mind is] focused}. %\footnote{The above two
%  paragraphs may still contain errors in readings. See footntoe
%  \fnref{note:avrtti} on page \pageref{note:avrtti}.}
%
At\refm{65} the stage of being focused, there is no domination by the stage.  Why?  Because defects (\emph{kleśa}) and \karman\ are not powerful [at that stage].
Because when the power of defects and \karman\ are powerful, [the states of] diversion and so on arise.   \emph{Samādhi\/} at the stage of being focused is \bht{the one that illuminates} = makes \bht{the object} understood \bht{as is} = as
itself. [The \Bhasyakara] uses the expression ``as is (\emph{%
  bhūta}),'' because, if there is even a trace of non-exactness (\emph{a\-yathābhūta}), the knowledge is not for a yogin.%
  \footnote{%
	There is some suspicion that sentence is corrupt. The first word \emph{ayogyarthajñānam} is already suspect. One would rather expect \emph{ayogyārtha°}. Also, the appearance of the word \emph{jñāna} is somewhat confusing. Knowledge has not been discussed in relation to \cittabhumi s and therefore appears out of context. One could also conceive that the compound was in fact two words, \emph{ayogy(a/ā)rthañ jñānam}; the conjunct character for \emph{ñjña} is complicated and can easily look like \emph{jña}. Or it is even possible that the word \emph{jñānam} is a corruption from something entirely different. Furthermore, the word \emph{gandhita} is not well attested in the sense our author apparently intends. \Sankara\ uses the word \emph{āgandhita} or its negated form \emph{anāgandhita} in the sense ``not left with any trace of,'' and this is one of his idiosyncrasies. We could consider the possibility that instead of \emph{ayathābhūtatvagandhitam}, the text in fact read \emph{ayathābhūtānāgandhitam}. Even after considering these emendations, the text does not appear very coherent.%
  }  [\emph{Samādhi\/} at the stage of being focused] \bht{destroys}---makes the five fold \bht{defects}, \ignorance, etc.,\footnote{See \rYS{2.03}{2.3} (\emph{avidyāsmitārāgadveṣābhiniveṣāḥ kleśāḥ}).} destroyed---
  \bht{[It] loosens the bindings of \karman}---[the compound \emph{karmabandhanāni\/} is a \KD\ compound in which] \karman s are the bindings [connected in the \emph{sāmānādhikaraṇya\/} relationship]. The [\karman s] are good, bad, and mixture of good and bad; or, [the compound \emph{karmabandhanāni\/} is a \TP\ compound that means] those originated from \karman s and they are the birth, etc., i.e., the bindings---\bht{[and it] loosens} = makes something loose. %
%[``Loosens'' means] to make something loose.\footnote{This sentence first
%  seems awkward since there is ``{yadi karmajāti} in the
%  interpretation.  This, however, becomes clearer after consulting
%  relevant sections of the PYS.  First, for the first member of the
%  compound {karma-bandhana}, which is being interpreted, {%
%    karman}, see YS \ldots, YBh \ldots.  For {bandhana}, there
%  is a mention in YBh i.e.\ on \rYS{1.24---}{}.  The three bondages is
%  explained as ``{prākṛta-vaikṛta-dākṣiṇāani
%    bandhanāni}'' by the author of the YVi while interpreting the PYS
%  1.24.  For the misleading compound ``{%
%    dharmādharmavimiśrāṇi}, it is believed that it should be taken
%  as dvandva compound ({ekaśeṣa}).  The author of the YVi
%  holds an opinion that karmans are of three kinds.  This is seen at
%  the interpretation on 1.5 (page ??), 4.7--8.  On the other hand,
%  according to the \rYS{4.07@4.7}, karmans of people but yogins are of three
%  kind.  Anyways as its being in the same case and number with {%
%    karmāṇi} it should be taken as explanation of the word.  With
%  regards to a following word {janmādīni}, it seems natural to
%  think that the author of the YVi is talking about the karmans of
%  people but yogins.  As the word {karmajāti} also appears in
%  the YBh 4.7, and there is another kind of karman the reading {%
%    yadi karmajāti} can be understood as it is talking about the
%  condition under which {vipākas} such as the birth ({%
%    janman}), and so on are produced.  The effects are birth, length of
%  life, experiences, according to the \rYS{4.9}.}
 \bht{[It] directs [the yogin] to suppression (\nirodha)}---orients him toward
[suppression].  \bht{It is called ``conscious (\emph{samprajñāta})''} by
the master(s).\footnote{\label{no_yogah}Note that the author of the YVi reads the \YBh\ \emph{\ldots\ nirodham āmukhīkaroti\dd\ sa samprajñātaḥ}, while most versions of the YBh reads \emph{sa samprajñataḥ yogaḥ}.  See \citet[3]{maas:2006}.  As a result of this lack of the word \emph{yogaḥ}, the author of the YVi considers the subject of the YBh to be \emph{samādhi\/} throughout after \emph{yogaḥ samādhiḥ}.  Hence, the terms \emph{samprajñāta\/} and \emph{asamprajñāta\/} can only modify the word \emph{samādhi}.  Consequently he considers the next \sutra\ as the definition of \emph{asamprajñāta-samādhi}.  See the next paragraph.}

The%\pratika{sa ca vitarkā°}
\pratikap{p01_70}\refm{66} \bhasya\ continues---\bht{And that is accompanied by \emph{vitarka}, or by \emph{vic\=ara}}, and so on.  [The reason why the \bhasya\ mentions this] is because \emph{samprajñāta\/} occurs first.  On the other hand, when it comes to stating the definition, since  \emph{asamprajñāta\/} is more important, it is reasonable to state the definition of it here---\bht{this} \emph{samprajñāta} \bht{we will teach later}.%\footnote{There could be other interpretation
%  for this section.  Coupled with difficulty in distinguishing the
%  text from the \bhasya\ and reconstructing the text of the YVi
%  itself, the sentence translated here might still have problems.
%  Consult the following paragraph where the author of the YVi exhibits
%  almost identical interpretation in slightlymore elaborately.  The
%  place where the \bhasya\ is referring to teach {%
%    samprajñāta} is \rYS{1.17}.}

For%\pratika{itaś cā°}
\pratikap{p01_71}\refm{87} the following reason, too, the definition of the \emph{%
  asamprajñāta\/} should be given at this very place.  Without
having the \emph{samprajñāta} prerequisite, the \emph{asamprajñāta} can
be realized by either excellent detachment or practicing the notion
of termination.\footnote{See \rYS{1.12--18}{1.12--18}.}  In order to show this [the
definition of the \emph{asamprajñāta} comes first].  On the other
hand, if the definition of \emph{samprajñāta\/}
was stated here, and if the definition of \emph{%
  asamprajñāta} was stated later, there would be a suspicion that one
qualifies for \emph{asamprajñāta-samādhi\/} with \emph{%
  samprajñāta\/} as prerequisite.  Therefore [the \Bhasyakara] states \bht{we will teach later}.

